\documentclass[12pt,a4paper]{article}
\usepackage[margin=2.2cm]{geometry}
\usepackage{amsmath,amssymb,mathtools}
\usepackage{booktabs,longtable,tabularx,array,multirow}
\usepackage{graphicx}
\usepackage{float}
\usepackage{xcolor}
\usepackage{hyperref}
\usepackage{url}
\usepackage{enumitem}
\usepackage{tikz}
\usetikzlibrary{arrows.meta,positioning,shapes.geometric,fit,calc}
\usepackage[most]{tcolorbox}
\usepackage{siunitx}
\usepackage{setspace}
\usepackage{caption}
\usepackage{amsmath,amssymb}
\usepackage{booktabs}



\usetikzlibrary{
	arrows.meta,
	angles,
	quotes,
	calc,
	positioning
}





% Persian typesetting; load xepersian after all other packages.
\usepackage{xepersian}
\settextfont{Amiri}
\setlatintextfont{DejaVu Serif}

\newcommand{\TikzFaText}[2][4cm]{%
	\parbox{#1}{%
		\setRTL
		\centering
		#2%
	}%
}



\linespread{1.28}
\hypersetup{colorlinks=true,linkcolor=blue!50!black,citecolor=green!40!black,urlcolor=blue!60!black}
\setlist[itemize]{itemsep=3pt,topsep=4pt}
\setlist[enumerate]{itemsep=3pt,topsep=4pt}

\definecolor{navy}{RGB}{25,55,95}
\definecolor{softblue}{RGB}{235,243,252}
\definecolor{softgreen}{RGB}{235,248,240}
\definecolor{softred}{RGB}{252,238,238}
\definecolor{softgray}{RGB}{245,245,245}

\newtcolorbox{keybox}[1][]{colback=softblue,colframe=navy,title=#1,fonttitle=\bfseries,breakable}
\newtcolorbox{warningbox}[1][]{colback=softred,colframe=red!60!black,title=#1,fonttitle=\bfseries,breakable}
\newtcolorbox{decisionbox}[1][]{colback=softgreen,colframe=green!45!black,title=#1,fonttitle=\bfseries,breakable}

\newcommand{\EN}[1]{\lr{#1}}
\newcommand{\degsec}{\ensuremath{^\circ/\mathrm{s}}}
\newcommand{\etal}{\EN{et al.}}

\title{\textbf{صورت‌بندی فنی و معماری پیشنهادی سامانهٔ بلادرنگ شناسایی، ردیابی و تثبیت هدف با دوربین تله و گیمبال \EN{Pan--Tilt}}\\[5pt]
\large طراحی برای پردازش روی رایانهٔ معمولی، سکوی متحرک و لرزان، و آزمون آزمایشگاهی با دیوار ویدئویی}
\author{گزارش فنی اولیه }
\date{تیر ۱۴۰۵ / ژوئیهٔ ۲۰۲۶}

\begin{document}
\maketitle

\begin{tcolorbox}[colback=softgray,colframe=black!55,breakable]
\textbf{دامنه و مرز کاربرد.}
این سند یک سامانهٔ الکترواپتیکیِ غیرمسلح برای پایش، تصویربرداری، شناسایی بصری و نگه‌داشتن هدف در مرکز تصویر را بررسی می‌کند. موضوع سند، کنترل خط دید دوربین و گیمبال است و شامل کنترل سلاح، درگیری، هدایت مهمات یا تصمیم‌گیری آسیب‌زا نیست. نمونه‌های صنعتی و عمومیِ سامانه‌های هوابرد صرفاً برای استخراج اصول معماری، تثبیت و آزمون ذکر شده‌اند.
\end{tcolorbox}

\tableofcontents
\newpage

\section{برداشت منظم از مسئله}

هدف پروژه ساخت سامانه‌ای است که یک دوربین تله با میدان دید باریک را روی گیمبال دو یا سه‌محوره نصب کند، هدف را در جریان ویدئو تشخیص دهد، موقعیت آن را در هر فریم تخمین بزند و به گیمبال فرمان دهد تا تصویر هدف همواره تا حد امکان در مرکز کادر باقی بماند. پردازش اصلی باید روی یک رایانهٔ متعارف انجام شود و سامانه در نهایت قابلیت نصب روی سکوی متحرک و لرزان، مانند کوادکوپتر یا وسیلهٔ هوابرد سریع، را داشته باشد. آزمون اولیه در محیط بسته و با نمایش ویدئوهای هدف روی دیوار ویدئویی انجام خواهد شد.

مسئله را باید به‌عنوان ترکیب پنج زیرمسئله دید، نه صرفاً طبقه‌بندی تصویر:

\begin{enumerate}
    \item \textbf{تثبیت خط دید:} حذف یا تضعیف اثر چرخش، لرزش و ارتعاش پایه با استفاده از ژیروسکوپ، انکودر و حلقهٔ داخلی گیمبال؛
    \item \textbf{آشکارسازی و شناسایی:} یافتن کادر هدف و تعیین کلاس یا مدل آن؛
    \item \textbf{ردیابی کوتاه‌مدت:} تخمین سریع موقعیت هدف بین دفعات اجرای آشکارساز؛
    \item \textbf{تخمین و پیش‌بینی حالت:} تخمین زاویه و نرخ زاویه‌ای خط دید، جبران تأخیر و ادامهٔ پیش‌بینی در انسداد یا افت موقت آشکارسازی؛
    \item \textbf{کنترل دیداری گیمبال:} تبدیل خطای پیکسلی به خطای زاویه‌ای و تولید فرمان نرم، محدود و پایدار برای محورهای \EN{pan} و \EN{tilt}.
\end{enumerate}








\subsection{صورت‌بندی سامانه به‌عنوان یک مسئله چندبخشی حلقه‌بسته}
\label{subsec:closed_loop_problem_formulation}

مسئله مورد بررسی را نمی‌توان صرفاً به‌صورت یک مسئله طبقه‌بندی تصویر
در نظر گرفت. طبقه‌بندی تنها مشخص می‌کند که محتوای مشاهده‌شده به کدام
رده از اهداف تعلق دارد، درحالی‌که در یک سامانه واقعی مبتنی بر دوربین
تله و سازوکار \lr{Pan--Tilt}، هدف باید ابتدا در میدان دید باریک دوربین
اخذ شود، در طول زمان ردیابی گردد، حرکت آینده آن پیش‌بینی شود و با وجود
لرزش‌ها و حرکت سکوی حامل، در مرکز تصویر باقی بماند. ازاین‌رو، مسئله
حاضر از پنج زیرمسئله به‌هم‌پیوسته تشکیل می‌شود:

\begin{enumerate}
	\item \textbf{آشکارسازی و اخذ اولیه هدف:}
	یافتن هدف در میدان دید محدود دوربین تله و تعیین موقعیت اولیه آن
	در صفحه تصویر؛
	
	\item \textbf{بازشناسی هدف:}
	تعیین نوع یا کلاس هدف و برآورد میزان اطمینان مدل نسبت به تصمیم
	اتخاذشده؛
	
	\item \textbf{ردیابی و پیش‌بینی حرکت:}
	تخمین وضعیت هدف در فریم‌های متوالی، پیش‌بینی جابه‌جایی آن و
	ادامه ردیابی در زمان افت یا قطع کوتاه‌مدت مشاهده؛
	
	\item \textbf{تثبیت خط دید و جبران اغتشاشات:}
	کاهش اثر لرزش، باد، ارتعاشات سازه‌ای، حرکت سکوی حامل و سایر
	اغتشاشاتی که باعث جابه‌جایی ناخواسته تصویر می‌شوند؛
	
	\item \textbf{کنترل حلقه‌بسته \lr{Pan--Tilt}:}
	تبدیل خطای مکانی هدف نسبت به مرکز تصویر به فرمان‌های حرکتی
	مناسب برای محورهای \lr{Pan} و \lr{Tilt}.
\end{enumerate}

شکل~\ref{fig:five_subproblems_pan_tilt} ارتباط این پنج زیرمسئله را
در قالب یک حلقه بسته ادراک--کنش نشان می‌دهد. در این ساختار، خروجی
بخش بینایی مستقیماً به کنترل‌کننده گیمبال ارسال می‌شود و حرکت گیمبال
نیز تصویر ورودی فریم بعدی را تغییر می‌دهد. بنابراین، عملکرد نهایی
سامانه نه فقط به دقت آشکارساز یا طبقه‌بند، بلکه به هماهنگی میان
پردازش تصویر، تخمین حالت، پیش‌بینی حرکت، تثبیت خط دید و کنترل
بلادرنگ وابسته است.

\begin{figure*}[t]
	\centering
	\includegraphics[
	width=0.97\textwidth,
	keepaspectratio
	]{figures/pan_tilt_five_subproblems.png}
	\caption{
		صورت‌بندی مسئله ردیابی هدف با دوربین تله و سامانه
		\lr{Pan--Tilt} به‌عنوان ترکیب پنج زیرمسئله به‌هم‌پیوسته.
		پس از آشکارسازی و اخذ اولیه هدف، نوع هدف بازشناسی شده و موقعیت
		آن در فریم‌های متوالی ردیابی و پیش‌بینی می‌شود. هم‌زمان، اثر
		لرزش‌ها و حرکت سکوی حامل بر خط دید جبران می‌گردد. کنترل‌کننده
		با استفاده از خطای موقعیت هدف نسبت به مرکز تصویر، فرمان‌های
		مناسب \lr{Pan} و \lr{Tilt} را تولید می‌کند. مسیر بازخورد پایین
		شکل نشان می‌دهد که سامانه یک زنجیره پردازش یک‌طرفه نیست، بلکه
		یک حلقه بسته بلادرنگ است که هدف آن حفظ پیوسته هدف در مرکز تصویر
		دوربین تله است.
	}
	\label{fig:five_subproblems_pan_tilt}
\end{figure*}





\begin{keybox}[نتیجهٔ اصلی صورت‌بندی]
سامانه یک \textbf{حلقهٔ کنترل دیداری چندنرخی} است. دوربین، آشکارساز، ردیاب، تخمین‌گر حالت، ژیروسکوپ و کنترل‌کنندهٔ موتور در نرخ‌های یکسان کار نمی‌کنند. طراحی موفق باید صریحاً این نرخ‌های متفاوت، تأخیر تصویر، زمان استنتاج و محدودیت سرعت/شتاب گیمبال را در نظر بگیرد.
\end{keybox}

\subsection{فرض‌های عملیاتی لازم}

برای آنکه صورت مسئله قابل آزمون و قابل اندازه‌گیری باشد، باید پاکت عملیاتی \EN{operational envelope} تعریف شود:

\begin{itemize}
    \item فاصلهٔ تقریبی هدف تا دوربین یا بازهٔ آن؛
    \item مؤلفهٔ سرعت نسبی عمود بر خط دید؛
    \item میدان دید افقی و عمودی در زوم مورد استفاده؛
    \item نرخ فریم، زمان نوردهی و رزولوشن؛
    \item بیشینهٔ سرعت و شتاب زاویه‌ای گیمبال؛
    \item دامنه و طیف فرکانسی ارتعاش پایه؛
    \item تأخیر کل از ورود نور به حسگر تا اعمال فرمان موتور؛
    \item اندازهٔ هدف در تصویر و کمینهٔ تعداد پیکسل قابل قبول برای تشخیص؛
    \item وضعیت آغاز کار: آیا هدف از ابتدا داخل تصویر است یا جهت اولیه از بیرون داده می‌شود؟
\end{itemize}

\begin{warningbox}[محدودیت اطلاعاتی دوربین تلهٔ تنها]
اگر هدف بیرون میدان دید باریک باشد و هیچ جهت اولیه، مختصات تقریبی، پیش‌بینی قبلی یا حسگر کمکی وجود نداشته باشد، هیچ الگوریتم بینایی نمی‌تواند وجود و جهت هدف را از تصویری که آن را نمی‌بیند استنتاج کند. بنابراین سامانهٔ تله‌محور می‌تواند \textbf{قفل را حفظ و بازیابی محلی کند}، اما تضمین اخذ اولیهٔ سراسری مستلزم حداقل یک سرنخ بیرونی است. در آزمون آزمایشگاهی می‌توان شروع هر سناریو را با حضور هدف در میدان دید یا با فرمان اولیهٔ اپراتور تعریف کرد.
\end{warningbox}

\section{چرا میدان دید باریک مسئله را دشوار می‌کند؟}

\subsection{نرخ زاویه‌ای خط دید}

سرعت خطی به‌تنهایی تعیین‌کنندهٔ دشواری نیست. اگر مؤلفهٔ سرعت نسبی عمود بر خط دید را با $v_{\perp}$ و فاصله را با $R$ نشان دهیم، تقریب مرتبهٔ اول نرخ زاویه‌ای خط دید برابر است با

\begin{equation}
\omega_{\mathrm{LOS}} \simeq \frac{v_{\perp}}{R}.
\label{eq:losrate}
\end{equation}

برای سرعت \SI{500}{\kilo\meter\per\hour}، یعنی تقریباً \SI{138.9}{\meter\per\second}، جدول \ref{tab:losrate} یک کران بدبینانه با فرض $v_{\perp}=v$ را نشان می‌دهد.

\begin{table}[H]
\centering
\caption{نرخ زاویه‌ای تقریبی خط دید برای سرعت نسبی عرضی \SI{500}{\kilo\meter\per\hour}}
\label{tab:losrate}
\begin{tabular}{ccc}
\toprule
فاصله $R$ & $\omega_{\mathrm{LOS}}$ بر حسب رادیان بر ثانیه & $\omega_{\mathrm{LOS}}$ بر حسب درجه بر ثانیه \\
\midrule
\SI{50}{\meter} & 2.778 & 159.2 \\
\SI{100}{\meter} & 1.389 & 79.6 \\
\SI{500}{\meter} & 0.278 & 15.9 \\
\SI{1000}{\meter} & 0.139 & 8.0 \\
\bottomrule
\end{tabular}
\end{table}

نتیجه روشن است: گیمبالی با سقف \SI{180}{\degree\per\second} ممکن است در فاصلهٔ \SI{50}{\meter} از نظر سرعت اسمی به مرز اشباع برسد، در حالی که همان هدف در فاصلهٔ \SI{500}{\meter} نرخ زاویه‌ای بسیار قابل‌کنترل‌تری دارد. علاوه بر سرعت، شتاب زاویه‌ای، اینرسی بار، تأخیر و محدودیت جریان موتور نیز تعیین‌کننده‌اند.




%==============================================================
% بسته‌های مورد نیاز در پیش‌گفتار
%==============================================================




%==============================================================
\subsubsection{نقش نرخ زاویه‌ای خط دید در دشواری ردیابی}
\label{subsubsec:los-angular-rate}
%==============================================================

سرعت خطی هدف به‌تنهایی معیار مناسبی برای تعیین دشواری ردیابی نیست.
آنچه مستقیماً مشخص می‌کند هدف با چه سرعتی در صفحه تصویر جابه‌جا شود،
\emph{نرخ زاویه‌ای خط دید}
یا
\lr{Line-of-Sight Angular Rate}
است. برای روشن‌شدن این موضوع، در
شکل~\ref{fig:los-geometry}،
مکان دوربین با \(O\)، مکان لحظه‌ای هدف با \(P(t)\)، فاصله دوربین
تا هدف با \(R\)، و زاویه خط دید با \(\theta\) نشان داده شده است.

\begin{figure}[t]
	\centering
	
	\resizebox{0.98\linewidth}{!}{%
		\begin{tikzpicture}[
			>=Latex,
			line cap=round,
			line join=round,
			every node/.style={font=\small},
			camera/.style={
				draw=black!75,
				fill=black!6,
				rounded corners=2pt,
				minimum width=1.55cm,
				minimum height=0.80cm,
				inner sep=3pt
			},
			target/.style={
				circle,
				draw=red!75!black,
				fill=red!12,
				minimum size=7mm,
				inner sep=0pt
			},
			explanation/.style={
				draw,
				rounded corners=2pt,
				inner sep=6pt,
				minimum height=1.25cm
			}
			]
			
			%----------------------------------------------------------
			% مختصات اصلی
			%----------------------------------------------------------
			\coordinate (O)  at (0,0);
			\coordinate (P)  at (7,2);
			\coordinate (Pp) at (6.45,3.92);
			
			%----------------------------------------------------------
			% میدان دید باریک دوربین
			%----------------------------------------------------------
			\fill[blue!6]
			(O)
			--
			(10:8.45)
			arc[
			start angle=10,
			end angle=22,
			radius=8.45
			]
			--
			cycle;
			
			\draw[
			gray!70,
			dashed,
			thick
			]
			(O) -- (10:8.45);
			
			\draw[
			gray!70,
			dashed,
			thick
			]
			(O) -- (22:8.45);
			
			% زاویه میدان دید
			\draw[
			blue!65!black,
			->,
			thick
			]
			(10:1.55)
			arc[
			start angle=10,
			end angle=22,
			radius=1.55
			];
			
			\node[
			blue!65!black,
			font=\small
			]
			at (16:1.98)
			{$\alpha$};
			
			%----------------------------------------------------------
			% دوربین
			%----------------------------------------------------------
			\node[
			camera,
			anchor=east
			]
			at (-0.15,0)
			{\TikzFaText[1.7cm]{دوربین تله}};
			
			\fill[black] (O) circle (2pt);
			
			\node[
			below=2mm
			]
			at (O)
			{$O$};
			
			%----------------------------------------------------------
			% مکان فعلی و مکان بعدی هدف
			%----------------------------------------------------------
			\node[
			target
			]
			(targetP)
			at (P)
			{};
			
			\node[
			below right=1mm
			]
			at (P)
			{$P(t)$};
			
			\node[
			target,
			dashed
			]
			(targetPp)
			at (Pp)
			{};
			
			\node[
			above right=1mm
			]
			at (Pp)
			{$P(t+\Delta t)$};
			
			%----------------------------------------------------------
			% خط دید فعلی
			%----------------------------------------------------------
			\draw[
			blue!70!black,
			very thick,
			->
			]
			(O)
			--
			(P)
			node[
			midway,
			below=3mm,
			fill=white,
			inner sep=1pt
			]
			{$R=\|\mathbf r\|$};
			
			% خط دید در لحظه بعد
			\draw[
			blue!55!black,
			dashed,
			thick,
			->
			]
			(O)
			--
			(Pp);
			
			%----------------------------------------------------------
			% تغییر زاویه خط دید
			%----------------------------------------------------------
			\pic[
			draw=orange!85!black,
			->,
			thick,
			angle radius=13mm,
			angle eccentricity=1.35,
			"$\Delta\theta$"
			]
			{angle=P--O--Pp};
			
			%----------------------------------------------------------
			% جابه‌جایی عرضی هدف
			%----------------------------------------------------------
			\draw[
			red!75!black,
			very thick,
			->
			]
			(P)
			--
			(Pp)
			node[
			midway,
			right=2mm,
			fill=white,
			inner sep=1pt
			]
			{$\mathbf v_{\perp}\Delta t$};
			
			%----------------------------------------------------------
			% مؤلفه شعاعی سرعت
			%----------------------------------------------------------
			\draw[
			green!45!black,
			thick,
			->
			]
			(P)
			--
			($(P)!-0.28!(O)$)
			node[
			pos=0.70,
			below right=1mm,
			fill=white,
			inner sep=1pt
			]
			{$v_{\parallel}$};
			
			%----------------------------------------------------------
			% علامت تقریبی عمودبودن مؤلفه عرضی بر خط دید
			%----------------------------------------------------------
			\draw[
			black!60,
			thick
			]
			($(P)+(-0.05,0.20)$)
			--
			++(-0.19,-0.05)
			--
			++(0.05,-0.19);
			
			%----------------------------------------------------------
			% توضیح مؤلفه عرضی
			%----------------------------------------------------------
			\node[
			explanation,
			draw=red!55!black,
			fill=red!4
			]
			at (10.65,3.55)
			{%
				\TikzFaText[4.35cm]{%
					مؤلفه سرعت عمود بر خط دید، یعنی
					\(v_{\perp}\)، سبب تغییر زاویه خط دید
					و حرکت هدف در صفحه تصویر می‌شود.
				}%
			};
			
			%----------------------------------------------------------
			% توضیح مؤلفه شعاعی
			%----------------------------------------------------------
			\node[
			explanation,
			draw=green!45!black,
			fill=green!4
			]
			at (10.65,1.35)
			{%
				\TikzFaText[4.35cm]{%
					مؤلفه شعاعی \(v_{\parallel}\)،
					در تقریب مرتبه اول، بیشتر فاصله هدف
					و اندازه ظاهری آن را تغییر می‌دهد.
				}%
			};
			
			%----------------------------------------------------------
			% نتیجه هندسی
			%----------------------------------------------------------
			\node[
			explanation,
			draw=blue!55!black,
			fill=blue!4
			]
			at (4.25,-1.65)
			{%
				\TikzFaText[5.5cm]{%
					برای یک سرعت عرضی ثابت، هرچه فاصله
					\(R\) کمتر باشد، تغییر زاویه خط دید
					و سرعت حرکت هدف در تصویر بیشتر خواهد بود.
				}%
			};
			
		\end{tikzpicture}%
	}
	
	\caption{
		هندسه نرخ زاویه‌ای خط دید در ردیابی با دوربین تله.
		سرعت نسبی هدف نسبت به دوربین به دو مؤلفه شعاعی
		\(v_{\parallel}\)
		و عرضی
		\(\mathbf v_{\perp}\)
		تجزیه می‌شود. مؤلفه عرضی موجب تغییر زاویه خط دید و
		جابه‌جایی هدف در صفحه تصویر می‌شود؛ درحالی‌که مؤلفه شعاعی،
		در تقریب مرتبه اول، بیشتر فاصله هدف و اندازه ظاهری آن را
		تغییر می‌دهد. ازاین‌رو، دو هدف با سرعت خطی یکسان، بسته به
		فاصله و جهت حرکت خود نسبت به دوربین، می‌توانند دشواری
		ردیابی کاملاً متفاوتی ایجاد کنند.
	}
	\label{fig:los-geometry}
\end{figure}

برای بیان ریاضی این موضوع، بردار مکان نسبی هدف نسبت به دوربین و
بردار سرعت نسبی آن را به‌ترتیب به‌صورت

\begin{equation}
	\mathbf r
	=
	\mathbf p_{\mathrm{target}}
	-
	\mathbf p_{\mathrm{camera}},
	\qquad
	\mathbf v_{\mathrm{rel}}
	=
	\dot{\mathbf p}_{\mathrm{target}}
	-
	\dot{\mathbf p}_{\mathrm{camera}}
	\label{eq:relative-position-velocity}
\end{equation}

تعریف می‌کنیم. فاصله هدف از دوربین و بردار واحد خط دید برابرند با

\begin{equation}
	R=\|\mathbf r\|,
	\qquad
	\widehat{\mathbf r}
	=
	\frac{\mathbf r}{R}.
	\label{eq:range-unit-los}
\end{equation}

سرعت نسبی هدف را می‌توان به مؤلفه موازی با خط دید و مؤلفه عمود بر
آن تجزیه کرد:

\begin{equation}
	\mathbf v_{\mathrm{rel}}
	=
	v_{\parallel}\widehat{\mathbf r}
	+
	\mathbf v_{\perp},
	\label{eq:relative-velocity-decomposition}
\end{equation}

که در آن

\begin{equation}
	v_{\parallel}
	=
	\widehat{\mathbf r}^{\mathsf T}
	\mathbf v_{\mathrm{rel}},
	\qquad
	\mathbf v_{\perp}
	=
	\left(
	\mathbf I
	-
	\widehat{\mathbf r}
	\widehat{\mathbf r}^{\mathsf T}
	\right)
	\mathbf v_{\mathrm{rel}}.
	\label{eq:radial-transverse-components}
\end{equation}

مؤلفه
\(v_{\parallel}\widehat{\mathbf r}\)
در امتداد خط دید قرار دارد و عمدتاً فاصله هدف را تغییر می‌دهد.
در مقابل،
\(\mathbf v_{\perp}\)
عمود بر خط دید است و عامل اصلی جابه‌جایی زاویه‌ای هدف در تصویر
محسوب می‌شود.

در فضای سه‌بعدی، اندازه نرخ زاویه‌ای خط دید از رابطه زیر به‌دست
می‌آید:

\begin{equation}
	\boxed{
		\left\|
		\boldsymbol{\omega}_{\mathrm{LOS}}
		\right\|
		=
		\frac{
			\left\|
			\mathbf r
			\times
			\mathbf v_{\mathrm{rel}}
			\right\|
		}{
			R^{2}
		}
		=
		\frac{
			\left\|
			\mathbf v_{\perp}
			\right\|
		}{
			R
		}
	}
	\label{eq:los-angular-rate}
\end{equation}

رابطه~\eqref{eq:los-angular-rate} نشان می‌دهد که نرخ زاویه‌ای خط
دید با مؤلفه سرعت عمود بر خط دید نسبت مستقیم و با فاصله هدف نسبت
معکوس دارد. بنابراین، افزایش سرعت خطی تنها زمانی ردیابی را دشوارتر
می‌کند که بخش قابل‌توجهی از این سرعت در راستای عمود بر خط دید قرار
گرفته باشد.

در حالت دوبعدی، اگر بردار مکان نسبی به‌صورت

\begin{equation}
	\mathbf r
	=
	\begin{bmatrix}
		x\\
		y
	\end{bmatrix},
	\qquad
	\theta
	=
	\operatorname{atan2}(y,x)
	\label{eq:two-dimensional-los-angle}
\end{equation}

تعریف شود، مشتق زاویه خط دید برابر است با

\begin{equation}
	\dot{\theta}
	=
	\frac{
		x\dot y-y\dot x
	}{
		x^{2}+y^{2}
	}.
	\label{eq:exact-two-dimensional-los-rate}
\end{equation}

از سوی دیگر، چون

\begin{equation}
	R=\sqrt{x^{2}+y^{2}},
\end{equation}

می‌توان اندازه نرخ زاویه‌ای را به‌شکل ساده‌تر زیر نوشت:

\begin{equation}
	\boxed{
		|\dot{\theta}|
		=
		\frac{
			|v_{\perp}|
		}{
			R
		}
	}
	\label{eq:simple-two-dimensional-los-rate}
\end{equation}

در یک بازه زمانی کوتاه
\(\Delta t\)،
اگر تغییرات
\(R\)
و
\(v_{\perp}\)
نسبتاً محدود باشند، تغییر زاویه خط دید تقریباً برابر است با

\begin{equation}
	\Delta\theta
	\simeq
	\frac{
		v_{\perp}
	}{
		R
	}
	\Delta t.
	\label{eq:short-time-angular-displacement}
\end{equation}

رابطه~\eqref{eq:short-time-angular-displacement} بیان می‌کند که یک
جابه‌جایی عرضی یکسان، در فاصله نزدیک‌تر، تغییر زاویه بزرگ‌تری ایجاد
می‌کند. برای مثال، اگر سرعت خطی هدف برابر با

\begin{equation}
	500~\mathrm{km/h}
	=
	138.9~\mathrm{m/s}
\end{equation}

باشد، دست‌کم دو حالت کاملاً متفاوت قابل تصور است.

\begin{itemize}
	\item
	اگر هدف تقریباً در امتداد خط دید به سمت دوربین یا در خلاف جهت
	آن حرکت کند، آنگاه
	\(v_{\perp}\simeq0\)
	است. در این حالت، با وجود سرعت خطی زیاد، مرکز هدف ممکن است
	به‌آرامی در تصویر جابه‌جا شود؛ هرچند فاصله و اندازه ظاهری آن
	می‌توانند به‌سرعت تغییر کنند.
	
	\item
	اگر همان هدف تقریباً عمود بر خط دید حرکت کند، در حالت بحرانی
	خواهیم داشت
	\(v_{\perp}\simeq138.9~\mathrm{m/s}\).
	در این وضعیت، هدف ممکن است با نرخ زاویه‌ای بسیار زیاد از میدان
	دید باریک دوربین تله عبور کند.
\end{itemize}

برای حالت بحرانی دوم، نرخ زاویه‌ای خط دید در چند فاصله نمونه در
جدول~\ref{tab:los-rate-example} ارائه شده است.

\begin{table}[t]
	\centering
	\caption{
		اثر فاصله بر نرخ زاویه‌ای خط دید برای سرعت عرضی
		\(v_{\perp}=138.9~\mathrm{m/s}\)
		و میدان دید افقی
		\(\alpha=2^{\circ}\).
		زمان عبور، مدت تقریبی لازم برای حرکت هدف از یک لبه میدان دید
		تا لبه مقابل است.
	}
	\label{tab:los-rate-example}
	
	\begin{tabular}{cccc}
		\toprule
		فاصله
		&
		نرخ زاویه‌ای
		&
		زمان عبور
		&
		تعداد فریم
		\\
		\(R~(\mathrm{m})\)
		&
		\(|\dot{\theta}|~(\mathrm{deg/s})\)
		&
		\(T_{\mathrm{cross}}~(\mathrm{ms})\)
		&
		در \(30~\mathrm{fps}\)
		\\
		\midrule
		\(100\)
		&
		\(79.6\)
		&
		\(25.1\)
		&
		\(0.75\)
		\\
		
		\(500\)
		&
		\(15.9\)
		&
		\(125.7\)
		&
		\(3.77\)
		\\
		
		\(1000\)
		&
		\(8.0\)
		&
		\(251.3\)
		&
		\(7.54\)
		\\
		\bottomrule
	\end{tabular}
\end{table}

اگر میدان دید افقی دوربین را با
\(\alpha\)،
برحسب رادیان، نشان دهیم، زمان تقریبی حرکت هدف از مرکز تصویر تا
لبه میدان دید برابر است با

\begin{equation}
	T_{\mathrm{center\rightarrow edge}}
	\simeq
	\frac{
		\alpha R
	}{
		2|v_{\perp}|
	}.
	\label{eq:center-to-edge-time}
\end{equation}

همچنین، زمان تقریبی عبور هدف از تمام عرض میدان دید از رابطه

\begin{equation}
	\boxed{
		T_{\mathrm{cross}}
		\simeq
		\frac{
			\alpha R
		}{
			|v_{\perp}|
		}
	}
	\label{eq:field-of-view-crossing-time}
\end{equation}

به‌دست می‌آید.

برای نمونه، در فاصله
\(R=100~\mathrm{m}\)،
با میدان دید افقی
\(\alpha=2^{\circ}\)
و سرعت عرضی
\(v_{\perp}=138.9~\mathrm{m/s}\)،
هدف تقریباً در مدت

\begin{equation}
	T_{\mathrm{cross}}
	\simeq
	25.1~\mathrm{ms}
\end{equation}

از تمام عرض تصویر عبور می‌کند. در مقابل، فاصله زمانی میان دو فریم
یک دوربین
\(30~\mathrm{fps}\)
برابر است با

\begin{equation}
	T_{\mathrm{frame}}
	=
	\frac{1}{30}
	\simeq
	33.3~\mathrm{ms}.
\end{equation}

بنابراین، در این مثال، هدف می‌تواند در زمانی کمتر از فاصله میان دو
فریم، کل میدان دید دوربین را طی کند. در چنین شرایطی، اجرای آشکارساز
روی فریم‌های متوالی، بدون استفاده از پیش‌بینی حرکت، برای حفظ قفل
کافی نخواهد بود.

نکته مهم دیگر آن است که در سکوی متحرک، نرخ حرکت مشاهده‌شده در
تصویر فقط ناشی از حرکت هدف نیست. چرخش دوربین، لرزش حامل، اغتشاشات
مکانیکی و خطای باقی‌مانده سامانه تثبیت نیز در حرکت ظاهری هدف نقش
دارند. یک مدل ساده برای نرخ زاویه‌ای مشاهده‌شده در تصویر عبارت است
از

\begin{equation}
	\dot{\theta}_{\mathrm{image}}
	\simeq
	\dot{\theta}_{\mathrm{LOS}}
	-
	\dot{\theta}_{\mathrm{camera}}
	+
	\dot{\theta}_{\mathrm{residual}},
	\label{eq:image-angular-rate}
\end{equation}

که در آن:

\begin{itemize}
	\item
	\(\dot{\theta}_{\mathrm{LOS}}\)
	نرخ واقعی تغییر خط دید ناشی از حرکت نسبی هدف و حامل است؛
	
	\item
	\(\dot{\theta}_{\mathrm{camera}}\)
	نرخ چرخش کنترل‌شده دوربین یا گیمبال است؛
	
	\item
	\(\dot{\theta}_{\mathrm{residual}}\)
	اثر باقی‌مانده لرزش، خطای تثبیت، انعطاف سازه، تأخیر کنترل و
	اغتشاشات مکانیکی را نشان می‌دهد.
\end{itemize}

ازاین‌رو، دشواری واقعی ردیابی باید بر اساس مجموعه‌ای از متغیرها،
شامل نرخ زاویه‌ای هدف در تصویر، میدان دید دوربین، نرخ فریم، زمان
نوردهی، تأخیر پردازش، سرعت و شتاب گیمبال و کیفیت تثبیت مکانیکی و
اینرسیایی ارزیابی شود؛ نه صرفاً بر مبنای سرعت خطی هدف.






\subsection{زمان عبور از میدان دید}

اگر میدان دید افقی را $\Phi$ بر حسب رادیان فرض کنیم، زمان تقریبی عبور هدف از تمام میدان دید در بدترین حرکت عرضی چنین است:


\begin{equation}
	T_{\mathrm{cross}}
	\simeq
	\frac{R\Phi}{v_{\perp}}.
	\label{eq:crossing}
\end{equation}

\begin{table}[H]
	\centering
	\caption{زمان عبور از میدان دید و تعداد فریم قابل مشاهده برای سرعت عرضی \lr{500 km/h}}
	\label{tab:crossing}
	
	\resizebox{\textwidth}{!}{%
		\begin{tabular}{cccccc}
			\toprule
			فاصله
			&
			میدان دید افقی
			&
			زمان عبور
			&
			فریم در \lr{30 fps}
			&
			فریم در \lr{60 fps}
			&
			فریم در \lr{120 fps}
			\\
			\midrule
			\lr{50 m}   & \lr{$2^\circ$} & \lr{12.6 ms}  & \lr{0.38} & \lr{0.75}  & \lr{1.51} \\
			\lr{50 m}   & \lr{$5^\circ$} & \lr{31.4 ms}  & \lr{0.94} & \lr{1.88}  & \lr{3.77} \\
			\lr{100 m}  & \lr{$2^\circ$} & \lr{25.1 ms}  & \lr{0.75} & \lr{1.51}  & \lr{3.02} \\
			\lr{100 m}  & \lr{$5^\circ$} & \lr{62.8 ms}  & \lr{1.88} & \lr{3.77}  & \lr{7.54} \\
			\lr{500 m}  & \lr{$2^\circ$} & \lr{125.7 ms} & \lr{3.77} & \lr{7.54}  & \lr{15.08} \\
			\lr{500 m}  & \lr{$5^\circ$} & \lr{314.2 ms} & \lr{9.42} & \lr{18.85} & \lr{37.70} \\
			\lr{1000 m} & \lr{$2^\circ$} & \lr{251.3 ms} & \lr{7.54} & \lr{15.08} & \lr{30.16} \\
			\bottomrule
		\end{tabular}%
	}
\end{table}




در فاصلهٔ \SI{100}{\meter} و میدان دید $2^\circ$، هدف در حدود \SI{25}{\milli\second} کل عرض تصویر را طی می‌کند؛ حتی نرخ \EN{120 fps} فقط حدود سه نمونه ایجاد می‌کند. این سناریو با آشکارساز کند یا حلقهٔ فرمان مبتنی بر موقعیت فریم قبلی پایدار نخواهد بود. راه‌حل، صرفاً انتخاب شبکهٔ عصبی بهتر نیست؛ باید نرخ فریم، پیش‌بینی، فرمان نرخ زاویه‌ای، تثبیت اینرسی و پاکت عملیاتی با هم طراحی شوند.

\subsection{تبدیل لرزش زاویه‌ای به جابه‌جایی پیکسلی}

در مرکز تصویر و برای زوایای کوچک، تقریب سادهٔ زیر مفید است:

\begin{equation}
\Delta u_{\mathrm{px}} \approx \frac{W}{\Phi_h}\,\Delta\theta,
\label{eq:jitterpixel}
\end{equation}

که $W$ عرض تصویر بر حسب پیکسل، $\Phi_h$ میدان دید افقی بر حسب درجه، و $\Delta\theta$ خطای زاویه‌ای بر حسب درجه است. برای تصویر \EN{1920 px} و میدان دید $2.9^\circ$، هر $0.004^\circ$ خطای زاویه‌ای تقریباً به $2.65$ پیکسل تبدیل می‌شود. ازاین‌رو حتی مشخصه‌های ظاهراً کوچک لرزش زاویه‌ای در انتهای تله قابل مشاهده‌اند.

محصول عمومی \EN{DJI Zenmuse H30} در باریک‌ترین حالت میدان دید قطری $2.9^\circ$ و برای گیمبال بازهٔ لرزش زاویه‌ای اعلام‌شدهٔ $\pm0.004^\circ$ در پرواز دارد \cite{djih30}. این مثال نشان می‌دهد که اپتیک تله و تثبیت باید به‌صورت یک سامانهٔ واحد تحلیل شوند؛ البته اعداد سازندگان تحت روش‌های آزمون یکسان اندازه‌گیری نشده‌اند و مقایسهٔ مستقیم آن‌ها باید با احتیاط انجام شود.

\subsection{تاری حرکت و زمان نوردهی}

اگر نرخ زاویه‌ای باقیمانده پس از تثبیت $\omega_r$ و فاصلهٔ کانونی معادل بر حسب پیکسل $f_x$ باشد، پهنای تقریبی تاری حرکت در راستای افقی چنین است:

\begin{equation}
b_{\mathrm{px}} \simeq f_x\,\omega_r\,t_e,
\label{eq:blur}
\end{equation}

که $t_e$ زمان نوردهی است. بنابراین افزایش نرخ فریم بدون کاهش نوردهی، لزوماً تاری را حل نمی‌کند. دوربین باید کنترل دستی نوردهی، بهره و ترجیحاً شاتر سراسری \EN{global shutter} داشته باشد. برای نمونه، یک دوربین صنعتی \EN{Basler ace 2} با حسگر شاتر سراسری، رزولوشن حدود پنج مگاپیکسل، نرخ تا حدود \EN{98.8 fps}، ورودی تحریک سخت‌افزاری و مانت \EN{C-mount} عرضه می‌شود \cite{basler98}. ذکر این نمونه به معنی الزام به همین مدل نیست؛ مشخصات مورد نیاز پروژه مهم‌تر از برند است.

\section{معماری کلان پیشنهادی}

\begin{figure}[H]
\centering
\begin{tikzpicture}[
node distance=7mm and 9mm,
block/.style={draw,rounded corners,align=center,minimum height=10mm,minimum width=31mm,fill=softblue},
fast/.style={draw,rounded corners,align=center,minimum height=10mm,minimum width=31mm,fill=softgreen},
ctrl/.style={draw,rounded corners,align=center,minimum height=10mm,minimum width=31mm,fill=softred},
line/.style={-Latex,thick}
]
\node[block] (cam) {دوربین تله\\فریم + زمان‌مهر};
\node[block,left=of cam] (imu) {\EN{IMU} پایه و گیمبال\\انکودر محورها};
\node[fast,below=of cam] (grab) {دریافت و پیش‌پردازش\\\EN{Frame Grabber}};
\node[block,below left=of grab] (detect) {آشکارساز سبک\\کادر + کلاس + اطمینان};
\node[fast,below right=of grab] (track) {ردیاب سریع\\هر فریم};
\node[fast,below=15mm of grab] (fusion) {همجوشی و تخمین حالت\\\EN{KF/EKF/UKF/IMM}};
\node[ctrl,below=of fusion] (super) {ماشین حالت\\اخذ، قفل، پیش‌بینی، بازیابی};
\node[ctrl,below=of super] (visual) {کنترل دیداری بیرونی\\فرمان نرخ \EN{pan/tilt}};
\node[ctrl,right=of visual] (inner) {حلقهٔ داخلی گیمبال\\ژیرو، موتور، انکودر};
\node[block,right=of inner] (plant) {پایهٔ متحرک و لرزان\\دوربین و هدف};

\draw[line] (cam) -- (grab);
\draw[line] (grab) -- (detect);
\draw[line] (grab) -- (track);
\draw[line] (detect) -- (fusion);
\draw[line] (track) -- (fusion);
\draw[line] (imu) |- (fusion);
\draw[line] (fusion) -- (super);
\draw[line] (super) -- (visual);
\draw[line] (visual) -- (inner);
\draw[line] (inner) -- (plant);
\draw[line] (plant.north) |- (cam.east);
\draw[line] (imu.east) -- (cam.west);
\end{tikzpicture}
\caption{معماری حلقه‌بستهٔ چندنرخی پیشنهادی}
\label{fig:architecture}
\end{figure}

\subsection{اصل حلقهٔ داخلی و حلقهٔ بیرونی}

ساختار مناسب، آبشاری است:

\begin{itemize}
    \item \textbf{حلقهٔ داخلی گیمبال:} روی کنترلر گیمبال اجرا می‌شود، از ژیروسکوپ و انکودر استفاده می‌کند و باید سریع‌تر از حلقهٔ بینایی باشد. وظیفهٔ آن دنبال‌کردن فرمان نرخ/زاویه و دفع اغتشاش پایه است. در یک گیمبال تجاری، به‌روزرسانی حسگر و تصحیح موتور تا \EN{2000 Hz} گزارش شده است \cite{gremsymanual}؛ این عدد نمونه‌ای از اختلاف مقیاس زمانی حلقهٔ مکانیکی و پردازش تصویر است.
    \item \textbf{حلقهٔ بیرونی دیداری:} روی رایانه اجرا می‌شود و خطای زاویه‌ای هدف، نرخ حرکت و پیش‌بینی تأخیر را به فرمان مرجع گیمبال تبدیل می‌کند.
\end{itemize}

اتصال مستقیم خطای پیکسلی به \EN{PWM} موتور یا ارسال فرمان موقعیت با نرخ پایین، برای میدان دید باریک و پایهٔ لرزان مناسب نیست. کنترلر رایانه باید فرمان نرخ زاویه‌ای یا مرجع هموارشده را به کنترلر داخلی گیمبال بدهد و وضعیت واقعی گیمبال را بازخوانی کند. استاندارد عمومی \EN{MAVLink Gimbal Protocol v2} نیز تفکیک \EN{Gimbal Manager} و \EN{Gimbal Device} و فرمان‌های وضعیت/جهت‌گیری را تعریف می‌کند \cite{mavlinkgimbal}.

\subsection{تقسیم نرخ‌های اجرایی}

جدول \ref{tab:rates} یک هدف مهندسی اولیه است؛ مقادیر نهایی باید با سخت‌افزار واقعی اندازه‌گیری شوند.

\begin{table}[H]
\centering
\caption{نرخ‌های پیشنهادی اولیه برای نمونهٔ آزمایشگاهی}
\label{tab:rates}
\small
\begin{tabularx}{0.97\textwidth}{>{\bfseries}p{3.2cm}p{2.7cm}X}
\toprule
ماژول & نرخ هدف & توضیح \\
\midrule
حلقهٔ داخلی گیمبال & \EN{500--2000 Hz} & تثبیت ژیروسکوپی و کنترل موتور؛ ترجیحاً مستقل از رایانه \\
نمونه‌برداری \EN{IMU} & \EN{200--1000 Hz} & ثبت حرکت پایه و خوراک پیش‌نگر \\
دریافت دوربین & \EN{60--120 fps} & شاتر سراسری، نوردهی کوتاه و زمان‌مهر دقیق \\
ردیاب کوتاه‌مدت & برابر نرخ دوربین & تخمین کادر در همهٔ فریم‌ها \\
آشکارساز/بازشناسی & \EN{5--30 Hz} & ناهمگام؛ بسته به توان پردازنده و اندازهٔ ورودی \\
تخمین‌گر حالت و کنترل بیرونی & \EN{60--200 Hz} & با پیش‌بینی بین فریم‌ها و استفاده از \EN{IMU} \\
ثبت داده & ناهمگام & بدون مسدودکردن مسیر بلادرنگ \\
\bottomrule
\end{tabularx}
\end{table}

\section{زنجیرهٔ ادراک: تشخیص، شناسایی و ردیابی}

\subsection{چرا طبقه‌بند تنها کافی نیست؟}

داده‌های فعلیِ تصویری برای یادگیری نوع هدف مفیدند، اما خروجی طبقه‌بند فقط نام کلاس است و محل هدف را در فریم مشخص نمی‌کند. کنترل \EN{pan--tilt} به مرکز کادر هدف نیاز دارد؛ بنابراین حداقل یک آشکارساز کادرمحور لازم است. راه عملی:

\begin{enumerate}
    \item پاک‌سازی نام‌ها و کلاس‌های دیتاست کامل محلی؛
    \item ساخت \EN{manifest} تصویر، کلاس، زیرکلاس و توضیح متناظر؛
    \item برچسب‌گذاری دستی یک زیرمجموعهٔ باکیفیت با \EN{bounding box}؛
    \item استفاده از همان تصاویر برای تولید توالی‌های مصنوعیِ کنترل‌شده و کادرهای دقیق؛
    \item آموزش آشکارساز سبک و در صورت نیاز یک طبقه‌بند ریزدانه روی برش هدف؛
    \item ارزیابی مستقل روی ویدئوی واقعی و نه صرفاً تصاویر مرجع.
\end{enumerate}

\subsection{ساختار پیشنهادی آشکارساز--ردیاب}

برای رایانهٔ معمولی، اجرای شبکهٔ آشکارساز در هر فریم الزاماً بهترین انتخاب نیست. معماری ناهمگام زیر مناسب‌تر است:

\begin{itemize}
    \item آشکارساز سبک در فریم‌های کلیدی یا هنگام کاهش اطمینان؛
    \item ردیاب بسیار سریع در همهٔ فریم‌ها؛
    \item تصحیح دوره‌ای ردیاب با خروجی آشکارساز برای جلوگیری از رانش؛
    \item پیش‌بینی کادر و ناحیهٔ جست‌وجو توسط فیلتر حالت؛
    \item اجرای آشکارساز ابتدا روی \EN{ROI} پیش‌بینی‌شده و سپس، در صورت شکست، روی کل فریم.
\end{itemize}

یک خط مبنای قابل دفاع شامل آشکارساز کوچک یک‌مرحله‌ای در قالب \EN{ONNX/OpenVINO} و دو ردیاب \EN{KCF} و \EN{CSRT} است. \EN{KCF} به‌عنوان یک فیلتر همبستگی بسیار سریع شناخته می‌شود \cite{kcf}؛ \EN{CSRT} با مدل اعتماد مکانی و کانالی، در برابر تغییر شکل و پس‌زمینه مقاوم‌تر ولی سنگین‌تر است \cite{csrt}. برای مسیر یادگیری‌شده و سبک، \EN{LightTrack} گزینهٔ مقایسه‌ای مناسب است و با هدف استقرار کم‌هزینه طراحی شده است \cite{lighttrack}.

\begin{decisionbox}[پیشنهاد اجرایی برای نسخهٔ نخست]
نسخهٔ نخست را با سه مسیر مقایسه‌ای بسازید: 
(۱) آشکارساز + \EN{KCF} برای سقف سرعت،
(۲) آشکارساز + \EN{CSRT} برای پایداری بیشتر، و
(۳) آشکارساز + ردیاب سیامی سبک برای نسخهٔ پژوهشی. انتخاب نهایی باید بر اساس تأخیر \EN{p95}، خطای مرکز و نرخ از دست‌دادن هدف روی همان رایانه انجام شود، نه بر اساس شهرت الگوریتم.
\end{decisionbox}

\subsection{نمایش اطمینان و تشخیص شکست}

ردیاب نباید بدون معیار کیفیت تا ابد به‌روزرسانی شود. خروجی هر فریم باید شامل این مقادیر باشد:

\begin{equation}
\mathcal{O}_k = \{B_k,\, c_k,\, q_k,\, t_k\},
\end{equation}

که $B_k$ کادر، $c_k$ کلاس، $q_k$ اطمینان و $t_k$ زمان‌مهر است. شکست احتمالی هنگامی اعلام می‌شود که ترکیبی از شرایط زیر رخ دهد:

\begin{itemize}
    \item افت پاسخ همبستگی یا اطمینان آشکارساز؛
    \item جهش ناممکن کادر نسبت به مدل حرکتی؛
    \item ناسازگاری کلاس در چند فریم متوالی؛
    \item خروج بخش عمدهٔ کادر از تصویر؛
    \item تفاوت زیاد اندازه/نسبت ابعاد با پیش‌بینی؛
    \item نوآوری غیرعادی در فیلتر حالت.
\end{itemize}

\section{هندسهٔ تصویر و تخمین حالت}

\subsection{تبدیل خطای پیکسلی به زاویه}

پس از کالیبراسیون دوربین و تعیین $(f_x,f_y,c_x,c_y)$، اگر مرکز کادر هدف $(u_k,v_k)$ باشد، خطای زاویه‌ای نسبت به محور اپتیکی به‌صورت زیر محاسبه می‌شود:

\begin{align}
 e_{\psi,k} &= \arctan\!\left(\frac{u_k-c_x}{f_x}\right),\\
 e_{\theta,k} &= \arctan\!\left(\frac{v_k-c_y}{f_y}\right).
\label{eq:pixelangle}
\end{align}

استفاده از خطای نرمال‌شدهٔ $2(u-c_x)/W$ بدون کالیبراسیون در زوم‌های مختلف، به تغییر بهرهٔ حلقه و رفتار ناپایدار می‌انجامد. برای لنز زوم، پارامترهای داخلی باید تابع وضعیت زوم باشند یا در چند نقطه کالیبره شوند.

\subsection{مدل حالت پیشنهادی}

برای هر محور می‌توان مدل شتاب تقریباً ثابت را به کار برد. حالت دوبعدی:

\begin{equation}
\mathbf{x}_k =
\begin{bmatrix}
\psi_k & \theta_k & \dot\psi_k & \dot\theta_k & \ddot\psi_k & \ddot\theta_k
\end{bmatrix}^{\!T}.
\end{equation}

مدل گسسته با فاصلهٔ زمانی $\Delta t$:

\begin{equation}
\mathbf{x}_{k+1}=\mathbf{F}(\Delta t)\mathbf{x}_k+\mathbf{G}\mathbf{w}_k,
\qquad
\mathbf{z}_k=\mathbf{H}\mathbf{x}_k+\mathbf{v}_k,
\end{equation}

که مشاهدهٔ دیداری $\mathbf{z}_k=[e_{\psi,k},e_{\theta,k}]^T$ است. اطلاعات ژیروسکوپ پایه، زاویه و نرخ گیمبال به‌عنوان ورودی معلوم یا مشاهدهٔ تکمیلی وارد می‌شوند.

برای حرکت نرم، فیلتر کالمن خطی کافی است؛ برای هندسه و نویز غیرخطی، \EN{EKF/UKF}؛ و برای مانورهای ناگهانی، یک \EN{IMM} با مدل‌های سرعت ثابت و شتاب‌دار مناسب‌تر است. مقاله‌ای در \EN{Aerospace} در سال ۲۰۲۵ یک زنجیرهٔ یکپارچهٔ آشکارساز سبک، \EN{CSRT} و فیلترهای \EN{EKF/UKF} را برای سکوی هوابرد دارای گیمبال بررسی کرده و بر اهمیت همجوشی دید و دادهٔ اینرسی در باد، لرزش و انسداد تأکید می‌کند \cite{kim2025}.

\subsection{جبران تأخیر}

تأخیر کل را چنین تفکیک می‌کنیم:

\begin{equation}
\tau_{\mathrm{tot}} =
\tau_{\mathrm{exposure}}+
\tau_{\mathrm{readout}}+
\tau_{\mathrm{transport}}+
\tau_{\mathrm{detect/track}}+
\tau_{\mathrm{control}}+
\tau_{\mathrm{actuator}}.
\end{equation}

فرمانی که بر اساس محل هدف در زمان $t-\tau_{\mathrm{tot}}$ محاسبه شود، در حرکت سریع عقب‌مانده است. بنابراین باید حالت به زمان اعمال فرمان پیش‌بینی شود:

\begin{equation}
\widehat{\mathbf{x}}(t+\tau_c)=
\Phi(\tau_{\mathrm{tot}}+\tau_c)\widehat{\mathbf{x}}(t),
\end{equation}

که $\tau_c$ افق کوتاه پیش‌بینی کنترل است. پژوهش منتشرشده در \EN{Automatica} نشان می‌دهد که تأخیر اندازه‌گیری تصویری و عدم قطعیت سینماتیکی از محدودیت‌های اصلی ردیابی روی سکوی تثبیت‌شده‌اند و استفاده از مشاهده‌گر/پیش‌بین تأخیر در کنترل اهمیت بنیادی دارد \cite{yang2022}.

\section{کنترل گیمبال و تثبیت خط دید}

\subsection{فرمان نرخ به‌جای فرمان موقعیت خام}

برای هر محور، یک قانون عملی اولیه می‌تواند چنین باشد:

\begin{equation}
\omega_{\mathrm{cmd}} =
\operatorname{sat}\!\left(
\widehat{\omega}_{\mathrm{LOS}} +
K_p e + K_d\dot e + K_i\int e\,dt,
\omega_{\max}
\right),
\label{eq:control}
\end{equation}

با محدودکنندهٔ شتاب و جرک:

\begin{align}
|\dot\omega_{\mathrm{cmd}}| &\le \alpha_{\max},\\
|\ddot\omega_{\mathrm{cmd}}| &\le j_{\max}.
\end{align}

ترم $\widehat{\omega}_{\mathrm{LOS}}$ خوراک پیش‌نگر از تخمین حرکت هدف است و ترم‌های بازخوردی خطای مرکز را اصلاح می‌کنند. فرمان باید در مختصات درست --- بدنه، زمین یا محور گیمبال --- تولید شود. استفادهٔ هم‌زمان از زاویه و نرخ پایه برای جبران حرکت کوادکوپتر ضروری است.

\subsection{دو نوع اغتشاش متفاوت}

باید بین این دو تمایز گذاشت:

\begin{enumerate}
    \item \textbf{حرکت کم‌فرکانس پایه:} رول، پیچ و یاو ناشی از مانور؛ با خوراک پیش‌نگر \EN{IMU} و گیمبال جبران می‌شود.
    \item \textbf{ارتعاش پرفرکانس:} موتور، ملخ، عدم بالانس، تشدید سازه و سوخت/جریان هوا؛ به جداساز مکانیکی، بالانس، طراحی سازه و حلقهٔ سریع نیاز دارد.
\end{enumerate}

فیلتر نرم‌افزاری تصویر جایگزین جداسازی مکانیکی نیست؛ زیرا تاری ایجادشده در زمان نوردهی را پس از ثبت نمی‌توان همیشه بازیابی کرد. نمونه‌های عمومی هوابرد مانند \EN{WESCAM MX-10} از تثبیت چهارمحوره، جداساز غیرفعال و \EN{IMU} نصب‌شده روی میز اپتیکی استفاده می‌کنند \cite{wescam}; این معماری نشان می‌دهد تثبیت خوب حاصل ترکیب مکانیک، اینرسی و پردازش است، نه یک الگوریتم واحد.

\subsection{ماشین حالت نظارتی}

\begin{figure}[H]
\centering
\begin{tikzpicture}[
state/.style={draw,rounded corners,minimum width=28mm,minimum height=10mm,align=center,fill=softblue},
line/.style={-Latex,thick},node distance=12mm and 14mm]
\node[state] (search) {جست‌وجوی محلی\\\EN{SEARCH}};
\node[state,right=of search] (acq) {اخذ هدف\\\EN{ACQUIRE}};
\node[state,right=of acq] (lock) {قفل پایدار\\\EN{LOCK}};
\node[state,below=of lock] (coast) {پیش‌بینی بدون مشاهده\\\EN{COAST}};
\node[state,left=of coast] (reacq) {بازیابی\\\EN{REACQUIRE}};
\node[state,left=of reacq] (safe) {حالت امن\\\EN{SAFE}};
\draw[line] (search) -- node[above]{کشف} (acq);
\draw[line] (acq) -- node[above]{تأیید چندفریمی} (lock);
\draw[line] (lock) -- node[right]{افت کوتاه} (coast);
\draw[line] (coast) -- node[below]{پیش‌بینی معتبر} (reacq);
\draw[line] (reacq) -- node[left]{کشف مجدد} (acq);
\draw[line] (reacq) -- node[below]{شکست/زمان بیشینه} (safe);
\draw[line,bend left=28] (safe) to node[above]{فرمان آغاز} (search);
\draw[line] (lock.north east) to[out=45,in=135,looseness=5] node[above]{اطمینان بالا} (lock.north west);
\end{tikzpicture}
\caption{ماشین حالت پیشنهادی برای جلوگیری از فرمان‌های ناپایدار هنگام افت ردیابی}
\label{fig:statemachine}
\end{figure}

در حالت \EN{COAST}، گیمبال برای چند ده تا چند صد میلی‌ثانیه بر اساس پیش‌بینی ادامه می‌دهد و مدل ظاهری ردیاب به‌روزرسانی نمی‌شود. در \EN{REACQUIRE}، ناحیهٔ جست‌وجو حول خط دید پیش‌بینی‌شده بزرگ می‌شود. اگر هدف بازیابی نشد، سامانه وارد حالت امن می‌شود و از حرکت نامحدود یا نوسان جلوگیری می‌کند.

\section{مشخصات سخت‌افزاری پیشنهادی برای نمونهٔ آزمایشگاهی}

\subsection{دوربین}

حداقل‌های پیشنهادی:

\begin{itemize}
    \item شاتر سراسری؛
    \item نرخ واقعی حداقل \EN{60 fps} و ترجیحاً \EN{100--120 fps} در \EN{ROI}؛
    \item کنترل دستی زمان نوردهی، بهره و گاما؛
    \item خروجی بدون فشرده‌سازی یا با تأخیر بسیار کم؛
    \item \EN{USB3 Vision}، \EN{GigE Vision} یا رابط مشابه پایدار؛
    \item ورودی تحریک سخت‌افزاری و امکان زمان‌مهر؛
    \item لنز \EN{C-mount} تله با اعوجاج و بازی مکانیکی کم؛
    \item امکان قفل فوکوس در فاز نخست.
\end{itemize}

دوربین‌های مصرفی با \EN{rolling shutter} و فشرده‌سازی داخلی ممکن است در لرزش و حرکت سریع اعوجاج هندسی و تأخیر متغیر ایجاد کنند. برای نمونهٔ تحقیقاتی، دوربین صنعتی با شاتر سراسری انتخاب مطمئن‌تری است.

\subsection{گیمبال}

ویژگی‌های لازم:

\begin{itemize}
    \item کنترل واقعی \EN{pan} و \EN{tilt} با فرمان نرخ؛
    \item بازخوانی زاویه، نرخ، وضعیت و خطا؛
    \item \EN{IMU} و انکودر داخلی؛
    \item سرعت و شتاب کافی بر اساس جدول \ref{tab:losrate}؛
    \item رابط مستند \EN{serial/CAN/Ethernet/MAVLink}؛
    \item امکان تنظیم بهره و فیلترهای لرزش؛
    \item ظرفیت بار با حاشیه برای دوربین و لنز تله؛
    \item اینرسی کم و بالانس دقیق در هر سه محور.
\end{itemize}

برای مقیاس مقایسه، \EN{Gremsy T7} سرعت کنترل‌شدهٔ بیشینهٔ \SI{180}{\degree\per\second} و بازهٔ لرزش زاویه‌ای اعلام‌شدهٔ $\pm0.02^\circ$ دارد \cite{gremsyt7}. این مشخصه نشان می‌دهد که برای سناریوهای نزدیک و سرعت عرضی بالا، حتی گیمبال سریع تجاری نیز ممکن است به حد خود برسد. مشخصهٔ لازم باید از رابطهٔ \eqref{eq:losrate} و حاشیهٔ شتاب تعیین شود.

\subsection{رایانه و رابط‌ها}

در نسخهٔ پایه، رایانهٔ \EN{x86-64} با پردازندهٔ چند‌هسته‌ای، \SI{16}{\giga\byte} حافظه و ذخیره‌ساز سریع کافی است؛ پردازندهٔ گرافیکی مجزا یک مزیت است اما نباید شرط عملکرد پایه باشد. نرم‌افزار باید دو حالت داشته باشد:

\begin{itemize}
    \item حالت \EN{CPU-only} با مدل کم‌حجم و \EN{ROI}؛
    \item حالت شتاب‌یافته با \EN{OpenVINO/ONNX Runtime GPU/CUDA} در صورت وجود.
\end{itemize}

مهم‌تر از نام پردازنده، اندازه‌گیری تأخیر انتهابه‌انتها و جلوگیری از صف‌شدن فریم‌هاست. سامانه باید همیشه جدیدترین فریم را پردازش کند و در ازدحام، فریم‌های قدیمی را کنار بگذارد.

\subsection{حسگرهای اینرسی و هم‌زمانی}

حداقل دو منبع وضعیت مفید است:

\begin{enumerate}
    \item \EN{IMU} پایه برای حرکت کوادکوپتر/سکو؛
    \item \EN{IMU} یا انکودر گیمبال برای وضعیت خط دید.
\end{enumerate}

همهٔ داده‌ها باید در یک ساعت مشترک یا با تخمین آفست زمانی ثبت شوند. حتی خطای زمانی چند میلی‌ثانیه در نرخ زاویه‌ای بالا به خطای قابل توجه زاویه تبدیل می‌شود:

\begin{equation}
\Delta\theta \simeq \omega\,\Delta t.
\end{equation}

برای نمونه، در \SI{80}{\degree\per\second}، خطای زمان‌مهر \SI{10}{\milli\second} برابر حدود $0.8^\circ$ است که در میدان دید $2^\circ$ بسیار بزرگ محسوب می‌شود.

\section{استفادهٔ صحیح از دیتاست موجود}

\subsection{نقش تصاویر چندگیگابایتی}

نبودن تصاویر کامل در مخزن عمومی مانع طراحی نیست؛ پردازش و ممیزی باید روی نسخهٔ محلی انجام شود. دادهٔ موجود سه نقش دارد:

\begin{enumerate}
    \item آموزش شناسایی کلاس کلی و مدل ریزدانه؛
    \item استخراج نمونه برای برچسب‌گذاری کادر و آموزش آشکارساز؛
    \item تولید ویدئوهای مصنوعی کنترل‌شده با حرکت، مقیاس، تاری و پس‌زمینهٔ متنوع.
\end{enumerate}

متن‌های توضیحی می‌توانند برای ساخت سلسله‌مراتب برچسب و کاهش خطای نام‌گذاری مفید باشند، اما هستهٔ نسخهٔ نخست باید آشکارسازی و ردیابی بلادرنگ باشد؛ یادگیری تصویر--متن یک افزونهٔ پژوهشی بعدی است.

\subsection{برچسب‌گذاری لازم}

برای شروع، به‌جای برچسب‌گذاری کل چندگیگابایت داده، یک زیرمجموعهٔ متوازن و متنوع انتخاب شود. معیار انتخاب:

\begin{itemize}
    \item هر کلاس و زیرکلاس مهم؛
    \item زوایای دید متفاوت؛
    \item اندازه‌های مختلف هدف؛
    \item پس‌زمینهٔ ساده و پیچیده؛
    \item کیفیت، نور و تاری متفاوت؛
    \item حذف یا خوشه‌بندی تصاویر تکراری و نزدیک به تکراری.
\end{itemize}

تقسیم آموزش/اعتبارسنجی/آزمون باید در سطح منبع یا خوشهٔ شباهت انجام شود، نه تصادفی در سطح فایل. تصاویر یک وسیله یا دنبالهٔ تقریباً یکسان نباید هم‌زمان در آموزش و آزمون قرار گیرند.

\subsection{تولید دادهٔ مصنوعی برای کنترل}

از تصاویر ثابت می‌توان توالی‌هایی با حقیقت مبنا تولید کرد:

\begin{equation}
B_k = \mathcal{T}(B_0; s_k,\varphi_k,x_k,y_k),
\end{equation}

که $\mathcal{T}$ تبدیل مقیاس، چرخش و انتقال است. سپس تخریب‌های زیر افزوده می‌شوند:

\begin{itemize}
    \item تاری حرکت با طول و جهت معلوم؛
    \item لرزش زاویه‌ای باندمحدود و چندفرکانسی؛
    \item کاهش وضوح و فشرده‌سازی؛
    \item مه، کنتراست کم و نور متغیر؛
    \item پوشیدگی جزئی و خروج موقت از تصویر؛
    \item تغییر سریع مقیاس و ظاهر؛
    \item تأخیر و حذف فریم.
\end{itemize}

در فاز پیشرفته، طیف لرزش مصنوعی با دادهٔ واقعی \EN{IMU} از کوادکوپتر جایگزین می‌شود. استاندارد \EN{MIL-STD-810H Method 514.8} تأکید می‌کند که پروفایل ارتعاش باید برای سکوی واقعی و محل نصب \textbf{tailor} شود و منحنی‌های عمومی جایگزین اندازه‌گیری محل نصب نیستند \cite{milstd810}. برای آزمون غیرنظامی تجهیزات نیز \EN{IEC 60068-2-64} چارچوب ارتعاش تصادفی پهن‌باند ارائه می‌کند \cite{iec60068}.

\section{طرح آزمون با دیوار ویدئویی}

\subsection{هدف آزمون}

دیوار ویدئویی برای اعتبارسنجی مرحله‌ای بسیار مناسب است، به شرط آنکه صرفاً یک ویدئوی نمایشی نباشد. سامانه باید بتواند مسیر، سرعت ظاهری، اندازهٔ هدف و زمان حقیقت مبنا را از مولد ویدئو دریافت کند. بدین ترتیب خطای مرکز و تأخیر واقعی اندازه‌گیری می‌شود.

\subsection{چیدمان پیشنهادی}

\begin{enumerate}
    \item دیوار ویدئویی یا پروژکتور با ابعاد و نرخ تازه‌سازی معلوم؛
    \item دوربین تله و گیمبال در فاصلهٔ معلوم؛
    \item نشانگرهای ثابت در گوشه‌های نمایشگر برای کالیبراسیون هندسی؛
    \item رایانهٔ مولد سناریو که موقعیت دقیق هدف را در هر فریم ثبت کند؛
    \item رایانهٔ پردازش یا فرآیند جداگانه برای سامانهٔ ردیابی؛
    \item در مرحلهٔ بعد، پایهٔ گیمبال روی میز حرکت دو یا سه‌محوره یا محرک لرزش نصب شود؛
    \item ثبت هم‌زمان ویدئو، فرمان، زاویهٔ گیمبال، \EN{IMU} و زمان‌ها.
\end{enumerate}

\subsection{نکتهٔ هندسی آزمون}

حرکت پیکسلی روی نمایشگر با حرکت زاویه‌ای یک هدف دور یکسان نیست، اما با دانستن فاصلهٔ دوربین تا نمایشگر و اندازهٔ فیزیکی مسیر می‌توان نرخ زاویه‌ای مؤثر را محاسبه کرد:

\begin{equation}
\theta(t)=\arctan\left(\frac{x_{\mathrm{screen}}(t)-x_0}{D}\right).
\end{equation}

برای شبیه‌سازی نرخ‌های زیاد، لازم است یا نمایشگر بزرگ و فاصله کم باشد یا خود پایهٔ گیمبال/دوربین نیز حرکت داده شود. نرخ تازه‌سازی نمایشگر نیز سقف سناریو را محدود می‌کند؛ نمایشگر \EN{60 Hz} نمی‌تواند حرکت نرم \EN{120 Hz} تولید کند.

\subsection{مراحل آزمون}

\begin{longtable}{p{1.1cm}p{4cm}p{9.2cm}}
\caption{مراحل پیشنهادی آزمون آزمایشگاهی}\label{tab:teststages}\\
\toprule
مرحله & وضعیت & هدف \\
\midrule
\endfirsthead
\toprule
مرحله & وضعیت & هدف \\
\midrule
\endhead
A & دوربین و گیمبال ثابت، هدف کند & صحت آشکارسازی، کالیبراسیون و علامت فرمان \\
B & هدف سریع، پایه ثابت & ارزیابی ردیاب، تأخیر و اشباع سرعت گیمبال \\
C & هدف سریع با افت/انسداد & ارزیابی \EN{COAST} و بازیابی \\
D & پایه با حرکت کم‌فرکانس & آزمون خوراک پیش‌نگر \EN{IMU} و تفکیک حرکت پایه از هدف \\
E & ارتعاش چندفرکانسی & سنجش تاری، خطای مرکز و نیاز جداساز مکانیکی \\
F & نور و فوکوس متغیر & آزمون نوردهی کوتاه، بهره و افت شناسایی \\
G & چند هدف در تصویر & انتخاب هدف، جلوگیری از تعویض هویت و قواعد مأموریت \\
H & اجرای طولانی & مصرف پردازنده، نشت حافظه، داغ‌شدن، افت فریم و پایداری \\
\bottomrule
\end{longtable}

\section{معیارهای ارزیابی}

\subsection{معیارهای ادراک}

\begin{itemize}
    \item \EN{mAP} آشکارسازی در اندازه‌های مختلف هدف؛
    \item دقت، فراخوان و \EN{F1} شناسایی کلاس؛
    \item موفقیت و دقت ردیابی، \EN{IoU} و خطای مرکز؛
    \item نرخ تعویض هویت در سناریوی چندهدفه؛
    \item نرخ شکست در تاری، انسداد، نور کم و حرکت سریع.
\end{itemize}

\subsection{معیارهای حلقهٔ کنترل}

خطای نرمال‌شدهٔ مرکز:

\begin{equation}
e_{n,k}=\sqrt{
\left(\frac{u_k-c_x}{W/2}\right)^2+
\left(\frac{v_k-c_y}{H/2}\right)^2}.
\end{equation}

معیارهای اصلی:

\begin{itemize}
    \item میانگین، میانه و صدک ۹۵ خطای مرکز بر حسب پیکسل و درجه؛
    \item درصد زمان حضور کامل یا جزئی هدف در میدان دید؛
    \item تعداد از دست‌رفتن قفل در هر دقیقه؛
    \item زمان بازیابی؛
    \item بیشینهٔ فراجهش، زمان نشست و نوسان پایدار؛
    \item \EN{IAE}: $\int |e(t)|dt$ و \EN{ITAE}: $\int t|e(t)|dt$؛
    \item درصد زمان اشباع سرعت/شتاب گیمبال.
\end{itemize}

\subsection{معیارهای بلادرنگ}

فقط میانگین زمان کافی نیست. باید گزارش شود:

\begin{itemize}
    \item تأخیر \EN{p50/p95/p99} از زمان‌مهر فریم تا صدور فرمان؛
    \item تأخیر تا مشاهدهٔ اثر فرمان در انکودر؛
    \item نرخ فریم ورودی، پردازش‌شده و حذف‌شده؛
    \item طول صف‌ها و سن فریم پردازش‌شده؛
    \item مصرف هسته‌های پردازنده، حافظه و دما؛
    \item لرزش زمانی \EN{jitter} حلقهٔ فرمان.
\end{itemize}

\begin{decisionbox}[آستانه‌های اولیهٔ پیشنهادی برای نمونهٔ آزمایشگاهی]
این مقادیر الزام نهایی نیستند و پس از تعیین پاکت عملیاتی اصلاح می‌شوند: تأخیر دید تا فرمان در صدک ۹۵ کمتر از \SI{50}{\milli\second}، خطای مرکز صدک ۹۵ کمتر از ۵ درصد عرض/ارتفاع تصویر در محدودهٔ طراحی، حضور هدف در میدان دید بیش از ۹۹ درصد زمان سناریو، و بازیابی محلی کمتر از \SI{0.5}{\second} زمانی که هدف در پنجرهٔ پیش‌بینی باقی مانده است.
\end{decisionbox}

\section{مطالعات و سامانه‌های مشابه و درس‌های آن‌ها}

\begin{longtable}{p{3.6cm}p{3.1cm}p{7.6cm}}
\caption{منابع منتخب و ارتباط آن‌ها با پروژه}\label{tab:literature}\\
\toprule
منبع & حوزه & درس قابل استفاده \\
\midrule
\endfirsthead
\toprule
منبع & حوزه & درس قابل استفاده \\
\midrule
\endhead
Yang \etal، \EN{Automatica 2022} \cite{yang2022} & کنترل دیداری با تأخیر & تأخیر تصویر و عدم قطعیت سینماتیکی باید در مدل و پیش‌بینی کنترل وارد شوند. \\
Hansen و Figueiredo، \EN{Drones 2024} \cite{hansen2024} & ردیابی فعال با گیمبال & نگه‌داشتن هدف در مرکز، تاری و خطای بازفکنی را کاهش می‌دهد؛ گیمبال بخشی از ادراک فعال است. \\
Huynh \etal، \EN{Actuators 2024} \cite{huynh2024} & سرووی دیداری زمان محدود & مشاهده‌گر اغتشاش و کنترل مقاوم برای حرکت غیرقابل‌پیش‌بینی و اغتشاش خارجی مفید است. \\
Yang \etal، \EN{JINT 2021} \cite{yang2021} & دوربین \EN{pan--tilt} روی هواگرد & طراحی مستقیم در صفحه تصویر می‌تواند خطای ناشی از بازسازی سه‌بعدی را کاهش دهد. \\
Kim \etal، \EN{Aerospace 2025} \cite{kim2025} & زنجیرهٔ یکپارچهٔ کم‌منبع & ترکیب آشکارساز، \EN{CSRT}، کنترل گیمبال و کالمن برای باد، انسداد و نویز بررسی شده است. \\
Henriques \etal، \EN{TPAMI 2015} \cite{kcf} & ردیابی سریع & فیلتر همبستگی برای مسیر سریع پایه مناسب است، ولی باید مکانیزم کشف شکست داشته باشد. \\
Luke\v{z}i\v{c} \etal، \EN{CVPR 2017} \cite{csrt} & ردیابی مقاوم‌تر & اعتماد مکانی/کانالی رانش را کاهش می‌دهد و گزینهٔ پایهٔ مقاوم روی \EN{CPU} است. \\
Yan \etal، \EN{CVPR 2021} \cite{lighttrack} & ردیاب یادگیری‌شدهٔ سبک & مسیر مناسبی برای مقایسه با ردیاب‌های کلاسیک در محدودیت محاسباتی است. \\
\EN{WESCAM MX-10} \cite{wescam} & مرجع صنعتی هوابرد & تثبیت چندمحوره، جداساز غیرفعال و \EN{IMU} روی میز اپتیکی به‌صورت یکپارچه استفاده می‌شوند. \\
\EN{DJI H30} \cite{djih30} & دوربین زوم صنعتی & میدان دید بسیار باریک، خطای زاویه‌ای کوچک را به چند پیکسل تبدیل می‌کند؛ زوم و تثبیت جداشدنی نیستند. \\
\EN{Gremsy T7} \cite{gremsyt7} & گیمبال تجاری سریع & سرعت \SI{180}{\degree\per\second} معیار خوبی برای مقایسه با نرخ خط دید محاسبه‌شده است. \\
\bottomrule
\end{longtable}

کار Hansen و Figueiredo به‌طور مستقیم نشان می‌دهد که ردیابی فعال با گیمبال برای کاهش تاری ناشی از حرکت و لرزش و نگه‌داشتن جسم در مرکز میدان دید مفید است \cite{hansen2024}. پژوهش \EN{Actuators} نیز هدف کنترل را رساندن تصویر هدف به مرکز با خطای ماندگار کم و پاسخ گذرای نرم، در حضور حرکت نامعلوم و اغتشاش، تعریف می‌کند \cite{huynh2024}. بنابراین تعریف پروژه حاضر با ادبیات علمی کنترل دیداری و ادراک فعال هم‌راستاست، ولی نوآوری آن باید در پاکت سرعت بالا، میدان دید باریک، پردازش رایانهٔ معمولی، دادهٔ اختصاصی و اعتبارسنجی سخت‌افزاردرحلقه شکل بگیرد.

\section{ساختار نرم‌افزار پیشنهادی}

\begin{latin}
\small
\begin{verbatim}
DrHarati_PTZ/
|-- configs/
|   |-- camera.yaml
|   |-- gimbal.yaml
|   |-- models.yaml
|   `-- scenarios.yaml
|-- data_tools/
|   |-- audit_dataset.py
|   |-- build_manifest.py
|   |-- deduplicate.py
|   |-- export_annotations.py
|   `-- synthesize_video_sequences.py
|-- vision/
|   |-- detector.py
|   |-- classifier.py
|   |-- tracker_kcf.py
|   |-- tracker_csrt.py
|   |-- tracker_lightweight.py
|   `-- confidence_monitor.py
|-- estimation/
|   |-- calibration.py
|   |-- kalman.py
|   |-- imm.py
|   `-- delay_predictor.py
|-- control/
|   |-- visual_servo.py
|   |-- rate_limiter.py
|   |-- state_machine.py
|   `-- gimbal_interface.py
|-- runtime/
|   |-- frame_grabber.py
|   |-- scheduler.py
|   |-- shared_state.py
|   |-- telemetry_logger.py
|   `-- watchdog.py
|-- evaluation/
|   |-- offline_metrics.py
|   |-- latency_report.py
|   |-- control_metrics.py
|   `-- plot_run.py
`-- tests/
\end{verbatim}
\end{latin}

نخ‌های اجرایی پیشنهادی:

\begin{itemize}
    \item نخ دریافت دوربین با بافر حلقوی یک یا دو فریم؛
    \item نخ آشکارساز ناهمگام؛
    \item نخ ردیاب و تخمین‌گر بلادرنگ؛
    \item نخ ارتباط گیمبال با اولویت بالاتر؛
    \item نخ ثبت داده با صف محدود و حذف کنترل‌شده؛
    \item \EN{watchdog} برای قطع فرمان در توقف نرم‌افزار یا دادهٔ منقضی.
\end{itemize}

\section{نقشهٔ اجرای پروژه}

\subsection{بستهٔ کاری صفر: تثبیت نیازمندی}

خروجی این مرحله باید یک جدول رسمی از بازهٔ فاصله، میدان دید، سرعت نسبی، اندازهٔ هدف، نرخ فریم، نور، ارتعاش و مشخصات گیمبال باشد. بدون این جدول، عبارت «ردیابی هدف با سرعت ۵۰۰ کیلومتر بر ساعت» آزمون‌پذیر نیست.

\subsection{بستهٔ کاری یک: نمونهٔ نرم‌افزاری بدون گیمبال}

\begin{itemize}
    \item ممیزی دیتاست کامل محلی؛
    \item برچسب‌گذاری کادر زیرمجموعه؛
    \item آموزش آشکارساز سبک؛
    \item مقایسهٔ سه ردیاب؛
    \item تولید سناریوهای مصنوعی و گزارش سرعت \EN{CPU-only}؛
    \item پیاده‌سازی فیلتر حالت و منطق شکست.
\end{itemize}

\subsection{بستهٔ کاری دو: حلقه‌بسته روی دیوار ویدئویی}

\begin{itemize}
    \item کالیبراسیون دوربین و گیمبال؛
    \item کنترل نرخ \EN{pan/tilt}؛
    \item ثبت تأخیر و خطای مرکز؛
    \item سناریوهای سرعت و شتاب زاویه‌ای افزایشی؛
    \item آزمون اشباع، تأخیر و بازیابی.
\end{itemize}

\subsection{بستهٔ کاری سه: سخت‌افزاردرحلقه و اغتشاش}

\begin{itemize}
    \item نصب پایه روی میز حرکتی یا محرک؛
    \item تزریق حرکت کم‌فرکانس و لرزش چندفرکانسی؛
    \item هم‌زمانی \EN{IMU} و تصویر؛
    \item مقایسهٔ کنترل صرفاً دیداری با کنترل دیداری--اینرسی؛
    \item اندازه‌گیری اثر جداساز مکانیکی و تنظیم حلقهٔ گیمبال.
\end{itemize}

\subsection{بستهٔ کاری چهار: آزمون روی سکوی متحرک غیرمسلح}

پس از عبور از معیارهای آزمایشگاهی، سامانه روی یک سکوی ایمن و غیرمسلح با پاکت سرعت محدود آغاز می‌شود و به‌تدریج محدودهٔ حرکت گسترش می‌یابد. پروفایل ارتعاش واقعی اندازه‌گیری و به آزمایشگاه بازگردانده می‌شود تا آزمون‌های تکرارپذیر ساخته شوند.

\section{ریسک‌های اصلی و راه کاهش آن‌ها}

\begin{longtable}{p{4cm}p{5.2cm}p{5cm}}
\caption{ریسک‌های فنی پروژه}\label{tab:risks}\\
\toprule
ریسک & پیامد & راه کاهش \\
\midrule
\endfirsthead
\toprule
ریسک & پیامد & راه کاهش \\
\midrule
\endhead
میدان دید بسیار باریک & خروج هدف در چند فریم & نرخ فریم بالا، پیش‌بینی، فرمان نرخ، تعریف فاصلهٔ حداقل \\
تأخیر متغیر نرم‌افزار & نوسان و عقب‌ماندگی & پردازش جدیدترین فریم، صف محدود، اندازه‌گیری \EN{p99} \\
اشباع گیمبال & از دست‌دادن قفل & محاسبهٔ پاکت، محدودکنندهٔ شتاب، کاهش زوم یا افزایش فاصله \\
تاری حرکت & افت آشکارسازی و ردیابی & شاتر سراسری، نوردهی کوتاه، تثبیت اینرسی و جداساز \\
دادهٔ صرفاً تصویری ثابت & ضعف در ویدئوی واقعی & برچسب کادر، ویدئوی مصنوعی کنترل‌شده و جمع‌آوری ویدئوی واقعی \\
نشت داده و تصاویر مشابه & دقت غیرواقعی & تقسیم بر اساس منبع/خوشه، هش ادراکی \\
رانش ردیاب & تعقیب پس‌زمینه & تصحیح با آشکارساز، معیار اطمینان و توقف به‌روزرسانی \\
عدم هم‌زمانی \EN{IMU}/camera & خوراک پیش‌نگر غلط & سخت‌افزار زمان‌مهر، تخمین آفست و آزمون ضربه‌ای \\
نمایشگر آزمایش با نرخ پایین & سناریوی غیرواقعی & ثبت محدودیت، استفاده از \EN{120 Hz} یا میز حرکتی \\
مقایسهٔ اشتباه مشخصات سازندگان & انتخاب سخت‌افزار نامناسب & آزمون مستقل با دوربین/بار واقعی و روش اندازه‌گیری یکسان \\
\bottomrule
\end{longtable}

\section{تعریف نسخهٔ حداقل قابل ارائه}

نسخهٔ حداقل ولی معتبر پروژه باید این توانایی‌ها را به‌صورت مستند نشان دهد:

\begin{enumerate}
    \item تشخیص و شناسایی کلاس‌های منتخب روی ویدئوی دیوار؛
    \item محاسبهٔ خطای زاویه‌ای کالیبره‌شده از مرکز هدف؛
    \item ردیابی در نرخ کامل دوربین با اجرای دوره‌ای آشکارساز؛
    \item ارسال فرمان نرخ به گیمبال و نگه‌داشتن هدف در مرکز؛
    \item تخمین حرکت و ادامهٔ کوتاه‌مدت در افت مشاهده؛
    \item بازیابی هدف در پنجرهٔ محلی؛
    \item ثبت هم‌زمان فریم، کادر، فرمان، انکودر، \EN{IMU} و تأخیر؛
    \item گزارش خطای مرکز، زمان حضور در میدان دید و صدک‌های تأخیر؛
    \item مقایسهٔ حداقل دو ردیاب و دو روش کنترل؛
    \item آزمون با حرکت پایه و لرزش مصنوعی.
\end{enumerate}

\begin{keybox}[جمع‌بندی نهایی]
برای پروژهٔ حاضر مسیر صحیح، «طبقه‌بند سبک + یک \EN{PID} ساده» نیست. معماری قابل اتکا عبارت است از:
\[
\boxed{
\begin{gathered}
\text{تثبیت اینرسی}
\rightarrow
\text{آشکارسازی ناهمگام}
\rightarrow
\text{ردیابی سریع}
\\[2mm]
\text{تخمین و پیش‌بینی تأخیر}
\rightarrow
\text{کنترل نرخ گیمبال}
\rightarrow
\text{منطق بازیابی و ایمنی}
\end{gathered}
}
\]
میدان دید باریک، حرکت سریع و پایهٔ لرزان باعث می‌شوند که زمان‌مهر، نرخ زاویه‌ای، تأخیر، نوردهی و اشباع گیمبال به همان اندازهٔ دقت شبکهٔ عصبی مهم باشند. نخستین تصمیم مهندسی باید تعیین فاصله و نرخ زاویه‌ای طراحی باشد؛ سپس دوربین، گیمبال و الگوریتم بر همان اساس انتخاب و در دیوار ویدئویی و سخت‌افزاردرحلقه آزموده شوند.
\end{keybox}






%======================================================================
\section{گذار از آزمون نمایشی به آزمایشگاه ماژولار سخت‌افزار در حلقه}
\label{sec:modular_eo_hil}
%======================================================================

آزمون مبتنی بر پخش یک ویدئوی ازپیش‌ضبط‌شده روی پرده و فرمان‌دادن
به یک سازوکار \lr{Pan--Tilt}، اگرچه می‌تواند برای بررسی اولیه
اتصال دوربین، نرم‌افزار و موتور مفید باشد، به‌تنهایی نماینده یک
آزمون حلقه‌بسته معتبر نیست. در چنین آزمایشی، حرکت واقعی گیمبال
تأثیری بر محتوای ویدئوی در حال پخش ندارد. بنابراین، حتی اگر
کنترل‌کننده فرمان نامناسبی تولید کند، صحنه ورودی الزاماً مطابق
این فرمان تغییر نخواهد کرد.

ساختار یک آزمون ساده مبتنی بر ویدئوی ثابت را می‌توان به‌شکل زیر
نمایش داد:

\begin{equation}
	\text{ویدئوی ازپیش‌ضبط‌شده}
	\longrightarrow
	\text{دوربین}
	\longrightarrow
	\text{ردیاب}
	\longrightarrow
	\text{فرمان گیمبال}.
	\label{eq:open_loop_video_test}
\end{equation}

این ساختار از دید سامانه‌ای یک مسیر عمدتاً یک‌طرفه است. در مقابل،
در یک آزمایش معتبر، وضعیت واقعی یا شبیه‌سازی‌شده گیمبال باید زاویه
دید دوربین و در نتیجه تصویر فریم بعدی را تغییر دهد. ساختار مطلوب
به‌صورت زیر است:

\begin{equation}
	\begin{aligned}
		\text{تصویر}_{k}
		&\longrightarrow
		\text{آشکارسازی و ردیابی}
		\longrightarrow
		\text{کنترل‌کننده}
		\longrightarrow
		\text{فرمان گیمبال}_{k},
		\\
		\text{فرمان گیمبال}_{k}
		&\longrightarrow
		\text{وضعیت جدید گیمبال}
		\longrightarrow
		\text{بازرندر یا مشاهده صحنه}
		\longrightarrow
		\text{تصویر}_{k+1}.
	\end{aligned}
	\label{eq:closed_loop_hil}
\end{equation}

در این حالت، خطای کنترل یا تأخیر پردازش مستقیماً در تصویر بعدی
ظاهر می‌شود. اگر فرمان دیر صادر شود، هدف به لبه تصویر نزدیک
می‌شود؛ اگر سرعت گیمبال کافی نباشد، هدف از میدان دید خارج می‌شود؛
و اگر پیش‌بینی حرکت نادرست باشد، مدت زمان بازیابی هدف افزایش
می‌یابد. بنابراین، معیارهای عملکرد سامانه به‌صورت واقعی در یک
حلقه ادراک--کنش سنجیده می‌شوند.

\begin{warningbox}[محدودیت آزمون مبتنی بر پرده و ویدئوی ثابت]
	پرده یا نمایشگر زمانی بخشی از یک آزمایش علمی معتبر محسوب می‌شود
	که تصویر نمایش‌داده‌شده بر اساس وضعیت لحظه‌ای گیمبال، دوربین،
	سکوی حامل و هدف دوباره تولید شود. پخش یک فیلم مستقل از حرکت
	گیمبال صرفاً ورودی تصویری ایجاد می‌کند و یک حلقه بسته
	الکترواپتیکی کامل به وجود نمی‌آورد.
\end{warningbox}

%----------------------------------------------------------------------
\subsection{تعریف آزمایشگاه الکترواپتیکی سخت‌افزار در حلقه}
\label{subsec:eo_hil_definition}
%----------------------------------------------------------------------

چارچوب مناسب برای این پروژه را می‌توان یک
\lr{Modular Electro-Optical Hardware-in-the-Loop Testbed}
نامید. در این چارچوب، صحنه، سکوی حامل، دوربین، پردازشگر،
کنترل‌کننده و گیمبال به‌صورت ماژول‌های مستقل تعریف می‌شوند.
هر ماژول می‌تواند واقعی، شبیه‌سازی‌شده یا ترکیبی از هر دو باشد.

برای نمونه، ممکن است هدف و کوادکوپتر در محیط شبیه‌سازی قرار
داشته باشند، ولی الگوریتم ردیابی روی رایانه واقعی اجرا شود.
در مرحله‌ای دیگر، گیمبال واقعی جایگزین گیمبال مجازی گردد و
زاویه‌های اندازه‌گیری‌شده توسط انکودر واقعی به شبیه‌ساز بازگردانده
شوند. در مرحله کامل‌تر، دوربین واقعی نیز تصویری را دریافت کند که
توسط مولد صحنه و متناسب با وضعیت گیمبال تولید شده است.

مجموعه اجزای سامانه را به‌صورت زیر در نظر می‌گیریم:

\begin{equation}
	\mathcal{S}
	=
	\left\{
	\mathcal{G},
	\mathcal{P},
	\mathcal{E},
	\mathcal{C},
	\mathcal{V},
	\mathcal{K},
	\mathcal{A}
	\right\},
	\label{eq:hil_component_set}
\end{equation}

که در آن:

\begin{itemize}
	\item
	\(\mathcal{G}\):
	مولد سناریو و حرکت هدف؛
	
	\item
	\(\mathcal{P}\):
	مدل یا نمونه واقعی سکوی حامل؛
	
	\item
	\(\mathcal{E}\):
	مولد صحنه و مدل الکترواپتیکی دوربین؛
	
	\item
	\(\mathcal{C}\):
	منبع فریم، شامل دوربین مجازی یا واقعی؛
	
	\item
	\(\mathcal{V}\):
	پردازش بینایی، آشکارسازی، بازشناسی و ردیابی؛
	
	\item
	\(\mathcal{K}\):
	تخمین‌گر حالت، پیش‌بین و کنترل‌کننده دیداری؛
	
	\item
	\(\mathcal{A}\):
	گیمبال، محرک‌ها، انکودر و حسگرهای اینرسیایی.
\end{itemize}

اصل مهم در طراحی آن است که برای هر مؤلفه
\(\mathcal{S}_{i}\)
دست‌کم دو پیاده‌سازی قابل تعویض وجود داشته باشد:

\begin{equation}
	\mathcal{S}_{i}
	\in
	\left\{
	\mathcal{S}_{i}^{\mathrm{virtual}},
	\mathcal{S}_{i}^{\mathrm{real}},
	\mathcal{S}_{i}^{\mathrm{replay}}
	\right\}.
	\label{eq:replaceable_modules}
\end{equation}

در نتیجه، کد ردیابی و کنترل نباید به نوع خاصی از دوربین یا
گیمبال وابسته باشد. کنترل‌کننده فقط یک رابط عمومی برای دریافت
فریم، خواندن وضعیت گیمبال و ارسال فرمان نرخ یا زاویه مشاهده
می‌کند.

%----------------------------------------------------------------------
\subsection{معماری حلقه‌بسته و قابلیت جایگزینی اجزا}
\label{subsec:hil_modular_architecture}
%----------------------------------------------------------------------

شکل~\ref{fig:modular_hil_architecture} معماری پیشنهادی آزمایشگاه
را نشان می‌دهد. مسیر رو به پایین، زنجیره تولید صحنه، مشاهده،
پردازش و فرمان را نمایش می‌دهد و مسیر بازخورد، اثر وضعیت واقعی
گیمبال بر تصویر بعدی را وارد حلقه می‌کند.

\begin{figure}[H]
	\centering
	\resizebox{0.96\textwidth}{!}{%
		\begin{tikzpicture}[
			node distance=7mm,
			>=Latex,
			block/.style={
				draw=navy,
				fill=softblue,
				rounded corners=2mm,
				minimum width=13.8cm,
				minimum height=1.35cm,
				align=center,
				inner sep=5pt
			},
			physical/.style={
				draw=green!50!black,
				fill=softgreen,
				rounded corners=2mm,
				minimum width=13.8cm,
				minimum height=1.35cm,
				align=center,
				inner sep=5pt
			},
			process/.style={
				draw=orange!70!black,
				fill=orange!6,
				rounded corners=2mm,
				minimum width=13.8cm,
				minimum height=1.35cm,
				align=center,
				inner sep=5pt
			},
			feedback/.style={
				->,
				very thick,
				blue!65!black,
				rounded corners=4pt
			}
			]
			
			\node[block] (scenario) {%
				\TikzFaText[12.8cm]{%
					\textbf{مولد سناریو و حقیقت مبنا}\\
					مسیر، سرعت، شتاب، مانور، فاصله، اندازه و وضعیت مرئی یا پوشیده هدف
			}};
			
			\node[block,below=of scenario] (platform) {%
				\TikzFaText[12.8cm]{%
					\textbf{دوقلوی دیجیتال سکوی حامل}\\
					دینامیک کوادکوپتر، باد، حرکت پایه، ارتعاش، تغییر مرکز جرم و اغتشاشات
			}};
			
			\node[block,below=of platform] (scene) {%
				\TikzFaText[12.8cm]{%
					\textbf{مولد صحنه و مدل الکترواپتیکی}\\
					میدان دید، زوم، فاصله کانونی، نوردهی، تاری حرکت، نویز، فشرده‌سازی و تأخیر
			}};
			
			\node[physical,below=of scene] (camera) {%
				\TikzFaText[12.8cm]{%
					\textbf{منبع فریم قابل تعویض}\\
					دوربین مجازی، دوربین واقعی، جریان شبکه، بازپخش داده یا دوربین روبه‌روی نمایشگر فعال
			}};
			
			\node[process,below=of camera] (perception) {%
				\TikzFaText[12.8cm]{%
					\textbf{رایانه پردازش و ادراک}\\
					آشکارساز، بازشناس، ردیاب، تخمین‌گر حالت، پیش‌بینی حرکت و مدیریت افت مشاهده
			}};
			
			\node[process,below=of perception] (controller) {%
				\TikzFaText[12.8cm]{%
					\textbf{کنترل‌کننده دیداری}\\
					تبدیل خطای مرکز تصویر و حالت پیش‌بینی‌شده هدف به فرمان نرخ یا زاویه
			}};
			
			\node[physical,below=of controller] (gimbal) {%
				\TikzFaText[12.8cm]{%
					\textbf{گیمبال قابل تعویض}\\
					گیمبال مجازی، گیمبال واقعی، موتور، انکودر، حسگر اینرسیایی و محدودیت‌های سرعت و شتاب
			}};
			
			\draw[->,very thick,navy]
			(scenario) -- (platform);
			
			\draw[->,very thick,navy]
			(platform) -- (scene);
			
			\draw[->,very thick,navy]
			(scene) -- (camera);
			
			\draw[->,very thick,navy]
			(camera) -- (perception);
			
			\draw[->,very thick,navy]
			(perception) -- (controller);
			
			\draw[->,very thick,navy]
			(controller) -- (gimbal);
			
			\draw[feedback]
			(gimbal.east)
			-- ++(1.0,0)
			|- (scene.east);
			
			\node[
			draw=blue!60!black,
			fill=white,
			rounded corners=2pt,
			align=center,
			right=9mm of perception,
			text width=3.3cm
			] {%
				\TikzFaText[3.0cm]{%
					بازخورد زاویه، نرخ، تأخیر، اشباع و وضعیت واقعی گیمبال
			}};
			
		\end{tikzpicture}%
	}
	
	\caption{
		معماری ماژولار آزمایشگاه الکترواپتیکی سخت‌افزار در حلقه.
		هر بلوک می‌تواند به‌صورت مجازی، واقعی یا بازپخش‌شده پیاده‌سازی
		شود. مسیر بازخورد سبب می‌شود وضعیت گیمبال بر صحنه و تصویر فریم
		بعد اثر بگذارد؛ ازاین‌رو، سامانه یک حلقه واقعی ادراک--کنش است.
	}
	\label{fig:modular_hil_architecture}
\end{figure}

برای آنکه حلقه شبیه‌سازی از دید اطلاعاتی معتبر باشد، فریم
\(k+1\) باید تابعی از وضعیت لحظه‌ای هدف، سکوی حامل و گیمبال باشد:

\begin{equation}
	I_{k+1}
	=
	\mathcal{R}
	\left(
	\mathbf{x}^{\mathrm{target}}_{k+1},
	\mathbf{x}^{\mathrm{platform}}_{k+1},
	\mathbf{q}^{\mathrm{gimbal}}_{k+1},
	\boldsymbol{\eta}_{k+1}
	\right),
	\label{eq:hil_rendering_condition}
\end{equation}

که در آن
\(\mathcal{R}\)
عملگر تولید یا ثبت تصویر،
\(\mathbf{x}^{\mathrm{target}}\)
حالت هدف،
\(\mathbf{x}^{\mathrm{platform}}\)
حالت حامل،
\(\mathbf{q}^{\mathrm{gimbal}}\)
وضعیت گیمبال و
\(\boldsymbol{\eta}\)
مجموعه اغتشاشات اپتیکی و زمانی است.

اگر تصویر مستقل از
\(\mathbf{q}^{\mathrm{gimbal}}\)
تولید شود، سامانه از دید کنترل حلقه‌بسته ناقص خواهد بود.

%----------------------------------------------------------------------
\subsection{سطوح تدریجی تحقق آزمایشگاه}
\label{subsec:hil_realization_levels}
%----------------------------------------------------------------------

پیاده‌سازی آزمایشگاه نباید از ابتدا به کامل‌ترین و پرهزینه‌ترین
حالت محدود شود. یک ساختار مرحله‌ای اجازه می‌دهد ابتدا الگوریتم‌ها
در محیط امن و تکرارپذیر ارزیابی شوند و سپس اجزای واقعی به‌تدریج
وارد حلقه شوند.

\begin{table}[H]
	\centering
	\caption{سطوح پیشنهادی توسعه آزمایشگاه ماژولار}
	\label{tab:hil_levels}
	\small
	\begin{tabularx}{\textwidth}{
			>{\centering\arraybackslash}p{1.1cm}
			>{\raggedleft\arraybackslash}p{3.2cm}
			>{\raggedleft\arraybackslash}X
			>{\raggedleft\arraybackslash}p{4.0cm}
		}
		\toprule
		سطح &
		عنوان &
		اجزای واقعی &
		هدف اصلی آزمون
		\\
		\midrule
		
		\lr{0} &
		نرم‌افزار در حلقه &
		فقط الگوریتم ردیابی و کنترل &
		صحت منطقی حلقه، سناریوها و معیارهای ارزیابی
		\\
		
		\lr{1} &
		پردازشگر در حلقه &
		رایانه نهایی پردازش &
		اندازه‌گیری زمان اجرا، مصرف منابع و تأخیر واقعی
		\\
		
		\lr{2} &
		گیمبال در حلقه &
		موتور، انکودر و کنترلر گیمبال &
		بررسی اشباع، لقی، تأخیر، سرعت، شتاب و بازخورد واقعی
		\\
		
		\lr{3} &
		دوربین در حلقه &
		دوربین و لنز واقعی &
		ورود نوردهی، فوکوس، تاری، نویز و محدودیت اپتیکی
		\\
		
		\lr{4} &
		کنترلر پرواز در حلقه &
		کنترلر پرواز واقعی &
		بررسی ارتباط با حسگرها، زمان‌بندی و وضعیت حامل
		\\
		
		\lr{5} &
		پایه حرکتی در حلقه &
		دوربین و گیمبال روی پایه لرزان &
		آزمون تثبیت در برابر ارتعاش و حرکت بازتولیدشده
		\\
		
		\bottomrule
	\end{tabularx}
\end{table}

در سطح پردازشگر در حلقه، داده تصویری مجازی است ولی تمام پردازش
روی رایانه‌ای انجام می‌شود که قرار است در نمونه نهایی استفاده شود.
ازاین‌رو، زمان‌های زیر باید مستقیماً اندازه‌گیری شوند:

\begin{equation}
	T_{\mathrm{total}}
	=
	T_{\mathrm{capture}}
	+
	T_{\mathrm{detect}}
	+
	T_{\mathrm{track}}
	+
	T_{\mathrm{estimate}}
	+
	T_{\mathrm{control}}
	+
	T_{\mathrm{communication}}.
	\label{eq:hil_latency_budget}
\end{equation}

در سطح گیمبال در حلقه، فرمان کنترل‌کننده به گیمبال واقعی ارسال
می‌شود، ولی صحنه و دوربین می‌توانند مجازی باقی بمانند. زاویه
اندازه‌گیری‌شده توسط انکودر گیمبال واقعی به شبیه‌ساز بازگردانده
می‌شود و دوربین مجازی بر اساس همان زاویه فریم بعد را تولید می‌کند.
این سطح امکان بررسی بسیاری از محدودیت‌های واقعی محرک را بدون
نیاز به پرواز فراهم می‌سازد.

در سطح دوربین در حلقه، تصویر شبیه‌سازی‌شده توسط یک نمایشگر یا
پروژکتور فعال ارائه می‌شود و دوربین واقعی آن را ثبت می‌کند.
در این حالت، تصویر نمایشگر باید با وضعیت لحظه‌ای گیمبال همگام
باشد. بنابراین، نمایشگر نقش یک
\lr{Optical Scene Simulator}
را دارد و نه یک پخش‌کننده فیلم ثابت.

%----------------------------------------------------------------------
\subsection{نقش صحیح نمایشگر یا دیوار ویدئویی}
\label{subsec:active_display_role}
%----------------------------------------------------------------------

استفاده از نمایشگر بزرگ یا دیوار ویدئویی همچنان می‌تواند مفید
باشد، مشروط بر آنکه بخشی از حلقه بسته باشد. در این معماری، وضعیت
انکودر گیمبال به شبیه‌ساز ارسال می‌شود، دوربین مجازی صحنه را از
زاویه متناظر رندر می‌کند و تصویر جدید روی نمایشگر ارائه می‌گردد.
سپس دوربین واقعی همان تصویر را ثبت می‌کند.

\begin{equation}
	\begin{aligned}
		\text{زاویه واقعی گیمبال}
		&\longrightarrow
		\text{وضعیت دوربین مجازی}
		\\
		&\longrightarrow
		\text{رندر بلادرنگ صحنه}
		\\
		&\longrightarrow
		\text{نمایشگر فعال}
		\\
		&\longrightarrow
		\text{دوربین واقعی}
		\\
		&\longrightarrow
		\text{پردازش و فرمان بعدی}.
	\end{aligned}
	\label{eq:active_display_hil}
\end{equation}

در این سطح باید تأخیرهای زیر نیز کالیبره و ثبت شوند:

\begin{itemize}
	\item تأخیر رندر صحنه؛
	\item تأخیر ارسال و نمایش فریم؛
	\item نرخ تازه‌سازی نمایشگر؛
	\item زمان نوردهی دوربین؛
	\item تأخیر دریافت فریم؛
	\item اختلاف زمانی میان وضعیت گیمبال و تصویر متناظر.
\end{itemize}

اگر مجموع این تأخیرها اندازه‌گیری نشود، ممکن است عملکرد ضعیف
کنترل‌کننده به اشتباه به الگوریتم ردیابی نسبت داده شود، درحالی‌که
منشأ اصلی خطا، تأخیر نمایشگر یا ناهمگامی زمانی بوده است.

%----------------------------------------------------------------------
\subsection{الزامات اعتبار علمی و عملیاتی}
\label{subsec:hil_validity_requirements}
%----------------------------------------------------------------------

برای آنکه آزمایشگاه پیشنهادی در گزارش نهایی و جلسه دفاع
قابل اتکا باشد، دست‌کم الزامات زیر باید رعایت شوند:

\begin{enumerate}
	\item
	تمام فریم‌ها، فرمان‌ها و وضعیت‌های گیمبال باید دارای
	زمان‌مهر مشترک باشند.
	
	\item
	حقیقت مبنای موقعیت هدف، زاویه گیمبال و وضعیت حامل باید
	در حالت شبیه‌سازی ذخیره شود.
	
	\item
	سناریوها باید قابل تکرار و دارای بذر تصادفی ثبت‌شده باشند.
	
	\item
	محدودیت سرعت، شتاب، لقی، اشباع و تأخیر گیمبال باید در
	حلقه حضور داشته باشد.
	
	\item
	افت فریم، تأخیر متغیر، پوشیدگی و گم‌شدن موقت هدف باید
	به‌صورت کنترل‌شده ایجاد شوند.
	
	\item
	نتایج باید علاوه بر ویدئوی نمایشی، به‌صورت فایل خام
	\lr{CSV}، نمودار، گزارش زمان‌بندی و لاگ ذخیره شوند.
	
	\item
	هر نتیجه باید مشخص کند کدام اجزا واقعی و کدام اجزا
	مجازی بوده‌اند.
	
	\item
	هیچ نتیجه شبیه‌سازی‌شده‌ای نباید به‌عنوان نتیجه سخت‌افزار
	واقعی گزارش شود.
\end{enumerate}

معیارهای اصلی ارزیابی می‌توانند شامل موارد زیر باشند:

\begin{equation}
	\mathcal{M}
	=
	\left\{
	E_{\mathrm{center}},
	P_{\mathrm{lock}},
	T_{\mathrm{reacquire}},
	T_{\mathrm{total}},
	P_{\mathrm{loss}},
	S_{\mathrm{gimbal}},
	U_{\mathrm{CPU}},
	U_{\mathrm{GPU}}
	\right\},
	\label{eq:hil_metrics_set}
\end{equation}

که به‌ترتیب بیانگر خطای مرکز تصویر، نسبت زمان حفظ قفل، زمان
بازیابی، تأخیر انتهابه‌انتها، احتمال خروج هدف از تصویر، میزان
اشباع گیمبال و مصرف منابع پردازشی هستند.

\begin{keybox}[خروجی مورد انتظار آزمایشگاه]
	خروجی نهایی این بخش نباید صرفاً یک فیلم از حرکت گیمبال باشد.
	آزمایشگاه باید بتواند برای هر سناریو، ویدئوی همگام، وضعیت هدف،
	وضعیت گیمبال، فرمان کنترل، خطای مرکز تصویر، زمان‌بندی پردازش،
	رخدادهای گم‌شدن و بازیابی و میزان مصرف منابع را به‌طور خودکار
	ثبت و گزارش کند.
\end{keybox}

%----------------------------------------------------------------------
\subsection{معماری نهایی پیشنهادی}
\label{subsec:final_hil_architecture_decision}
%----------------------------------------------------------------------

\begin{decisionbox}[تصمیم طراحی]
	پروژه باید به‌صورت یک دوقلوی دیجیتال ماژولار آغاز شود و سپس
	رایانه پردازش، گیمبال، دوربین، کنترلر پرواز و پایه لرزان به‌صورت
	تدریجی وارد حلقه شوند. در این ساختار، نمایشگر یک جزء اختیاری
	برای آزمون اپتیکی دوربین است و نه هسته اصلی پروژه. هسته اصلی،
	بسته‌شدن حلقه میان وضعیت گیمبال، تصویر مشاهده‌شده، الگوریتم
	ردیابی و فرمان کنترلی است.
\end{decisionbox}

بنابراین، مسیر توسعه مناسب را می‌توان به‌صورت زیر خلاصه کرد:

\begin{equation}
	\boxed{
		\text{دوقلوی مناسب را می‌توان به‌صورت زیر خلاصه کرد:}
		\mathrel{\longrightarrow}
		\text{پردازشگر در حلقه}
		\mathrel{\longrightarrow}
		\text{گیمبال در حلقه}
		\mathrel{\longrightarrow}
		\text{دوربین در حلقه}
		\mathrel{\longrightarrow}
		\text{کنترلر پرواز و پایه حرکتی در حلقه}
	}
	\label{eq:hil_development_path}
\end{equation}


	
	در بخش بعد، زیرساخت‌ها، کتابخانه‌ها و پروژه‌های متن‌باز قابل
	استفاده برای تحقق این معماری بررسی می‌شوند.



\input{../latex_generated/hil_operational_results}









%======================================================================
\subsection{پیاده‌سازی و ارزیابی دموی گرافیکی حلقه‌بسته}
\label{subsec:graphical_closed_loop_demo}
%======================================================================

به‌منظور عبور از ارزیابی صرفاً عددی و فراهم‌کردن یک نمونه قابل
مشاهده برای ارزیابان صنعتی، یک دموی گرافیکی حلقه‌بسته برای سامانه
ردیابی تله‌فوتو و کنترل \lr{Pan--Tilt} توسعه داده شد. این دمو در
محیط ویندوز و با استفاده از فناوری‌های
\lr{WebGL/Three.js}
اجرا می‌شود و به‌صورت هم‌زمان سه بخش اصلی سامانه را نمایش می‌دهد:

\begin{enumerate}
	\item
	نمای سه‌بعدی محیط شامل سکوی حامل، گیمبال، دوربین و هدف متحرک؛
	
	\item
	تصویر دوربین تله‌فوتو شامل کادر هدف، مرکز تصویر، خطای افقی و
	عمودی و وضعیت جاری ردیابی؛
	
	\item
	داشبورد زنده شامل زاویه‌های \lr{Pan} و \lr{Tilt}، فرمان‌های
	کنترلی، خطای مرکز، وضعیت قفل و اطلاعات زمانی سامانه.
\end{enumerate}

برخلاف آزمون مبتنی بر پخش یک ویدئوی مستقل، در این دمو فرمان
کنترل‌کننده وضعیت گیمبال را تغییر می‌دهد و وضعیت جدید گیمبال نیز
جهت دوربین و تصویر بعدی را تعیین می‌کند. بنابراین، حلقه زیر به‌صورت
قابل مشاهده بسته می‌شود:

\begin{equation}
	I_k
	\longrightarrow
	\widehat{\mathbf z}_k
	\longrightarrow
	\mathbf u_k
	\longrightarrow
	\mathbf q^{\mathrm{gimbal}}_{k+1}
	\longrightarrow
	I_{k+1},
	\label{eq:implemented_demo_loop}
\end{equation}

که در آن
\(I_k\)
تصویر دوربین در فریم \(k\)،
\(\widehat{\mathbf z}_k\)
موقعیت تخمینی هدف در تصویر،
\(\mathbf u_k\)
فرمان کنترلی و
\(\mathbf q^{\mathrm{gimbal}}_{k+1}\)
وضعیت جدید گیمبال است.

\begin{keybox}[اجرای زنده دموی عملیاتی]
	دموی گرافیکی با یکی از فرمان‌های زیر از ریشه پروژه اجرا می‌شود:
	
	\begin{latin}
		\begin{verbatim}
			.\run_demo.ps1
		\end{verbatim}
		
		یا:
		
		\begin{verbatim}
			run_demo.bat
		\end{verbatim}
	\end{latin}
	
	پس از اجرا، برنامه در مرورگر باز شده و نمای سه‌بعدی، تصویر
	تله‌فوتو و داشبورد زنده را به‌طور هم‌زمان نمایش می‌دهد.
\end{keybox}

%----------------------------------------------------------------------
\subsubsection{نمای عملیاتی برنامه}
\label{subsubsec:operational_demo_view}
%----------------------------------------------------------------------

شکل~\ref{fig:operational_demo_main} نمای اصلی سامانه در حال اجرا
را نشان می‌دهد. حضور هم‌زمان محیط سه‌بعدی، مدل گیمبال، تصویر
تله‌فوتو، کادر هدف و مقادیر زنده کنترل، امکان مشاهده مستقیم ارتباط
میان ادراک و عمل را فراهم می‌کند.

\begin{figure}[H]
	\centering
	\includegraphics[
	width=0.96\textwidth,
	keepaspectratio
	]{figures/03_target_locked.png}
	
	\caption{
		نمای عملیاتی دموی گرافیکی حلقه‌بسته در حالت
		\lr{LOCK}.
		در این تصویر، مدل سه‌بعدی محیط، دوربین تله‌فوتو، کادر هدف،
		مرکز تصویر، وضعیت گیمبال و داده‌های زنده کنترل به‌طور هم‌زمان
		نمایش داده شده‌اند. این آزمایش در سطح نرم‌افزار در حلقه
		\lr{SIL}
		و با شبیه‌ساز گرافیکی مبتنی بر
		\lr{WebGL/Three.js}
		اجرا شده است.
	}
	\label{fig:operational_demo_main}
\end{figure}

%----------------------------------------------------------------------
\subsubsection{آزمون خاموش و روشن‌شدن کنترل‌کننده}
\label{subsubsec:controller_activation_test}
%----------------------------------------------------------------------

برای اثبات اثر واقعی کنترل‌کننده بر موقعیت هدف در تصویر، ابتدا
کنترل‌کننده غیرفعال شد. در این وضعیت، هدف به‌تدریج از مرکز تصویر
فاصله گرفت و خطای مرکز افزایش یافت. سپس کنترل‌کننده فعال شد؛
فرمان‌های \lr{Pan} و \lr{Tilt} تغییر کردند، مدل گیمبال چرخید و
هدف مجدداً به ناحیه مرکزی تصویر بازگردانده شد.

\begin{figure}[H]
	\centering
	
	\begin{minipage}[t]{0.48\textwidth}
		\centering
		\includegraphics[
		width=\linewidth,
		keepaspectratio
		]{figures/01_controller_off.png}
		
		\vspace{2mm}
		\TikzFaText[0.94\linewidth]{%
			\textbf{الف) کنترل‌کننده غیرفعال:}
			هدف از مرکز تصویر فاصله گرفته و خطای ردیابی افزایش یافته است.
		}
	\end{minipage}
	\hfill
	\begin{minipage}[t]{0.48\textwidth}
		\centering
		\includegraphics[
		width=\linewidth,
		keepaspectratio
		]{figures/02_gimbal_rotating.png}
		
		\vspace{2mm}
		\TikzFaText[0.94\linewidth]{%
			\textbf{ب) کنترل‌کننده فعال:}
			گیمبال در حال چرخش است و خطای مرکز به‌تدریج کاهش می‌یابد.
		}
	\end{minipage}
	
	\caption{
		مقایسه رفتار سامانه پیش و پس از فعال‌سازی کنترل‌کننده.
		در حالت غیرفعال، هدف بدون جبران از مرکز تصویر دور می‌شود. پس از
		فعال‌شدن کنترل‌کننده، فرمان‌های حرکتی موجب چرخش گیمبال، تغییر
		تصویر دوربین و بازگرداندن هدف به ناحیه مرکزی می‌شوند.
	}
	\label{fig:controller_off_on_comparison}
\end{figure}

نتایج ثبت‌شده در این آزمایش نشان می‌دهند که میانگین خطای مرکز در
حالت خاموش‌بودن کنترل‌کننده برابر با

\begin{equation}
	\overline{e}_{\mathrm{off}}
	=
	0.4941
\end{equation}

و در حالت فعال‌بودن آن برابر با

\begin{equation}
	\overline{e}_{\mathrm{on}}
	=
	0.00935
\end{equation}

بوده است. کاهش مطلق ثبت‌شده در خطای مرکز برابر با

\begin{equation}
	\Delta \overline{e}
	=
	\overline{e}_{\mathrm{off}}
	-
	\overline{e}_{\mathrm{on}}
	=
	0.4847
\end{equation}

گزارش شد. همچنین، در طول اجرای دمو، دامنه تغییر زاویه محورهای
گیمبال به‌ترتیب برابر با

\begin{equation}
	\Delta\psi_{\mathrm{Pan}}
	=
	6.97^\circ,
	\qquad
	\Delta\theta_{\mathrm{Tilt}}
	=
	0.74^\circ
	\label{eq:measured_gimbal_motion}
\end{equation}

اندازه‌گیری شد. این تغییرات نشان می‌دهند که کاهش خطا فقط نتیجه
تغییر عددی متغیرهای نرم‌افزاری نبوده و با تغییر قابل مشاهده جهت
گیمبال همراه شده است.

%----------------------------------------------------------------------
\subsubsection{اغتشاش، افت مشاهده و بازیابی هدف}
\label{subsubsec:demo_disturbance_recovery}
%----------------------------------------------------------------------

برای ارزیابی رفتار سامانه در شرایط غیرایده‌آل، اغتشاش لرزشی و
افت موقت مشاهده هدف نیز در سناریو وارد شد. در زمان فعال‌شدن
لرزش، تصویر و خط دید تحت اغتشاش قرار گرفتند، ولی کنترل‌کننده تلاش
کرد خطای باقی‌مانده را کاهش دهد.

همچنین، هنگام قطع موقت مشاهده، ماشین حالت از
\lr{LOCK}
به
\lr{COAST}
منتقل شد. در این وضعیت، سامانه بدون دریافت مشاهده جدید، حرکت هدف
را به‌طور موقت پیش‌بینی کرد. پس از بازگشت مشاهده، حالت
\lr{REACQUIRE}
فعال شد و در نهایت سامانه مجدداً به وضعیت
\lr{LOCK}
بازگشت.

\begin{figure}[H]
	\centering
	
	\begin{minipage}[t]{0.48\textwidth}
		\centering
		\includegraphics[
		width=\linewidth,
		keepaspectratio
		]{figures/04_vibration_enabled.png}
		
		\vspace{2mm}
		\TikzFaText[0.94\linewidth]{%
			\textbf{الف) فعال‌بودن لرزش:}
			اثر اغتشاش حامل بر تصویر و تلاش کنترل‌کننده برای جبران آن.
		}
	\end{minipage}
	\hfill
	\begin{minipage}[t]{0.48\textwidth}
		\centering
		\includegraphics[
		width=\linewidth,
		keepaspectratio
		]{figures/05_coast.png}
		
		\vspace{2mm}
		\TikzFaText[0.94\linewidth]{%
			\textbf{ب) حالت \lr{COAST}:}
			ادامه تخمین حرکت در زمان قطع موقت مشاهده هدف.
		}
	\end{minipage}
	
	\vspace{5mm}
	
	\begin{minipage}[t]{0.48\textwidth}
		\centering
		\includegraphics[
		width=\linewidth,
		keepaspectratio
		]{figures/06_reacquire.png}
		
		\vspace{2mm}
		\TikzFaText[0.94\linewidth]{%
			\textbf{ج) حالت \lr{REACQUIRE}:}
			جست‌وجوی محلی و بازیابی هدف پس از بازگشت مشاهده.
		}
	\end{minipage}
	\hfill
	\begin{minipage}[t]{0.48\textwidth}
		\centering
		\includegraphics[
		width=\linewidth,
		keepaspectratio
		]{figures/03_target_locked.png}
		
		\vspace{2mm}
		\TikzFaText[0.94\linewidth]{%
			\textbf{د) بازگشت به \lr{LOCK}:}
			قرارگرفتن مجدد هدف در نزدیکی مرکز تصویر.
		}
	\end{minipage}
	
	\caption{
		مراحل ارزیابی مقاومت سامانه در برابر لرزش و قطع موقت مشاهده.
		توالی حالت‌های ثبت‌شده شامل
		\lr{LOCK}
		به
		\lr{COAST}،
		سپس
		\lr{REACQUIRE}
		و در نهایت بازگشت به
		\lr{LOCK}
		بوده است.
	}
	\label{fig:demo_state_transitions}
\end{figure}

توالی حالات ثبت‌شده در فایل رخدادهای آزمایش به‌صورت زیر است:

\begin{equation}
	\boxed{
		\mathrm{LOCK}
		\longrightarrow
		\mathrm{COAST}
		\longrightarrow
		\mathrm{REACQUIRE}
		\longrightarrow
		\mathrm{LOCK}
	}
	\label{eq:measured_state_transition}
\end{equation}

%----------------------------------------------------------------------
\subsubsection{خلاصه خروجی‌های قابل بازتولید}
\label{subsubsec:demo_reproducible_outputs}
%----------------------------------------------------------------------

جدول~\ref{tab:graphical_demo_results} مهم‌ترین مشخصات و خروجی‌های
ثبت‌شده دمو را خلاصه می‌کند.

\begin{table}[H]
	\centering
	\caption{خلاصه نتایج ثبت‌شده دموی گرافیکی حلقه‌بسته}
	\label{tab:graphical_demo_results}
	\small
	\begin{tabularx}{\textwidth}{
			>{\raggedleft\arraybackslash}p{5.5cm}
			>{\centering\arraybackslash}X
		}
		\toprule
		شاخص یا خروجی &
		مقدار ثبت‌شده
		\\
		\midrule
		
		نوع آزمایش &
		\lr{Graphical Software-in-the-Loop}
		\\
		
		زیرساخت گرافیکی &
		\lr{Windows-native WebGL/Three.js}
		\\
		
		تفکیک ویدئوی دفاعی &
		\lr{1280 $\times$ 800}
		\\
		
		مدت ویدئوی دفاعی &
		\lr{35.0 s}
		\\
		
		تعداد نمونه‌های تله‌متری &
		\lr{3504}
		\\
		
		میانگین خطا با کنترل‌کننده خاموش &
		\lr{0.4941}
		\\
		
		میانگین خطا با کنترل‌کننده روشن &
		\lr{0.00935}
		\\
		
		کاهش مطلق خطای مرکز &
		\lr{0.4847}
		\\
		
		دامنه حرکت محور \lr{Pan} &
		\lr{6.97 deg}
		\\
		
		دامنه حرکت محور \lr{Tilt} &
		\lr{0.74 deg}
		\\
		
		توالی حالات مشاهده‌شده &
		\lr{LOCK $\rightarrow$ COAST $\rightarrow$ REACQUIRE
			$\rightarrow$ LOCK}
		\\
		
		وجود مقادیر نامعتبر &
		\lr{No NaN/Inf}
		\\
		
		وضعیت اجرای Webots &
		اجرا نشده است
		\\
		
		وضعیت اجرای Gazebo/ArduPilot &
		اجرا نشده است؛ مربوط به فاز دوم
		\\
		
		\bottomrule
	\end{tabularx}
\end{table}

خروجی‌های اصلی برای بازتولید و بررسی آزمایش عبارت‌اند از:

\begin{itemize}
	\item
	فایل اجرای زنده:
	\path{run_demo.ps1}
	و
	\path{run_demo.bat}؛
	
	\item
	ویدئوی کامل دمو:
	\path{results/videos/defense_demo.mp4}؛
	
	\item
	داده تله‌متری:
	\path{results/raw/telemetry.csv}؛
	
	\item
	رخدادها و تغییر حالات:
	\path{results/raw/events.json}؛
	
	\item
	خلاصه معیارها:
	\path{results/summary/demo_metrics.json}؛
	
	\item
	تصاویر مراحل مختلف آزمایش:
	\path{results/screenshots/}.
\end{itemize}

\begin{keybox}[شاهد اصلی عملیاتی‌بودن حلقه]
	شاهد اصلی این آزمایش صرفاً کاهش یک متغیر عددی نیست. پس از فعال‌شدن
	کنترل‌کننده، فرمان‌های \lr{Pan/Tilt} تغییر می‌کنند، مدل گیمبال
	به‌طور قابل مشاهده می‌چرخد، جهت دوربین و تصویر بعدی تغییر می‌کنند
	و هدف به مرکز تصویر بازمی‌گردد. این توالی، وجود بازخورد میان ادراک،
	کنترل و تصویر را در سطح نرم‌افزار در حلقه نشان می‌دهد.
\end{keybox}

%----------------------------------------------------------------------
\subsubsection{سطح اعتبار و محدودیت‌های نسخه فعلی}
\label{subsubsec:demo_current_limitations}
%----------------------------------------------------------------------

با وجود عملیاتی‌بودن دموی گرافیکی، سطح فعلی نباید با یک سامانه
\lr{Hardware-in-the-Loop}
واقعی یکسان تلقی شود. اجزای زیر در این مرحله به‌صورت مجازی
پیاده‌سازی شده‌اند:

\begin{itemize}
	\item سکوی حامل و دینامیک آن؛
	\item گیمبال و محرک‌های آن؛
	\item دوربین و مدل تصویربرداری؛
	\item هدف و حرکت آن؛
	\item اغتشاشات و افت مشاهده.
\end{itemize}

همچنین، در این مرحله
\lr{Webots}،
\lr{Gazebo}،
\lr{ArduPilot SITL}،
کنترلر پرواز واقعی، گیمبال واقعی، دوربین واقعی و ارتباط
\lr{MAVLink}
اجرا نشده‌اند. بنابراین، طبقه‌بندی صحیح نتیجه فعلی عبارت است از:

\begin{equation}
	\boxed{
		\text{\lr{Interactive Graphical SIL Demonstrator}}
	}
	\label{eq:demo_correct_classification}
\end{equation}

و نه یک نمونه کامل
\lr{EO-HIL}.

نکته مهم دیگر آن است که در گزارش نهایی باید منشأ کادر هدف و
اندازه‌گیری تصویری صریحاً مشخص شود. اگر موقعیت کادر مستقیماً از
حقیقت مبنای شبیه‌ساز دریافت شده باشد، این خروجی نباید به‌عنوان
نتیجه آشکارساز یا ردیاب تصویری معرفی شود. در نسخه کامل، لازم است
چهار کمیت زیر به‌صورت مستقل ثبت و قابل مقایسه باشند:

\begin{equation}
	\left\{
	\mathbf z_k^{\mathrm{ground\ truth}},
	\mathbf z_k^{\mathrm{detector}},
	\mathbf z_k^{\mathrm{tracker}},
	\widehat{\mathbf z}_{k|k-1}^{\mathrm{prediction}}
	\right\}.
	\label{eq:perception_sources_separation}
\end{equation}

در صورتی که کادر فعلی از پردازش واقعی پیکسل‌ها استخراج شده باشد،
نام الگوریتم، نرخ اجرا، تأخیر و confidence آن نیز باید در جدول
نتایج ذکر شود.

\begin{warningbox}[محدودیت تفسیر نتایج]
	اعداد گزارش‌شده مربوط به یک سناریوی مشخص از دموی گرافیکی هستند و
	نباید بدون اجرای چند سناریو، چند بذر تصادفی و شرایط متفاوت سرعت،
	فاصله، میدان دید، تأخیر و لرزش به‌عنوان عملکرد عمومی سامانه
	تفسیر شوند. این نتایج اثبات می‌کنند که حلقه گرافیکی و منطق کنترل
	فعال است، ولی هنوز جایگزین اعتبارسنجی چندسناریویی و آزمون
	سخت‌افزاری نیستند.
\end{warningbox}

%----------------------------------------------------------------------
\subsubsection{گام بعدی توسعه}
\label{subsubsec:demo_next_stage}
%----------------------------------------------------------------------

این دمو پایه مناسبی برای نمایش دفاع و بررسی منطق حلقه‌بسته فراهم
کرده است. گام بعدی باید بدون تغییر رابط‌های نرم‌افزاری موجود،
اجزای مجازی را به‌ترتیب با اجزای استاندارد یا واقعی جایگزین کند:

\begin{equation}
	\begin{aligned}
		\text{\lr{WebGL/Three.js SIL}}
		&\longrightarrow
		\text{\lr{Image-based detector/tracker}}
		\\
		&\longrightarrow
		\text{\lr{Gazebo + ArduPilot SITL}}
		\\
		&\longrightarrow
		\text{\lr{Real gimbal in the loop}}
		\\
		&\longrightarrow
		\text{\lr{Real camera in the loop}}.
	\end{aligned}
	\label{eq:demo_to_hil_path}
\end{equation}

\begin{decisionbox}[نتیجه مرحله فعلی]
	نسخه فعلی برای اثبات دیداری مفهوم، نمایش حلقه کنترل و ارائه در
	جلسه دفاع مناسب است؛ زیرا ارزیاب می‌تواند حرکت هدف، چرخش گیمبال،
	تغییر تصویر، کاهش خطا و حالات بازیابی را مستقیماً مشاهده کند.
	بااین‌حال، ارزش اصلی این مرحله در فراهم‌کردن یک پایه اجرایی
	قابل تعویض است و مرحله بعد باید بر اتصال همان نرم‌افزار به
	شبیه‌ساز استاندارد پرواز و سپس سخت‌افزار واقعی متمرکز شود.
\end{decisionbox}








%======================================================================
\subsection{اجرای حلقه‌بسته استاندارد مبتنی بر Gazebo و ArduPilot}
\label{subsec:gazebo_ardupilot_closed_loop}
%======================================================================

پس از اعتبارسنجی اولیه منطق کنترل در دموی گرافیکی مرورگر، نسخه‌ای
استانداردتر از سامانه با استفاده از
\EN{Gazebo Harmonic}
و
\EN{ArduPilot SITL}
پیاده‌سازی و اجرا شد. هدف از این مرحله آن بود که حلقه ادراک--کنترل
دیگر صرفاً بر یک محیط نمایشی اختصاصی متکی نباشد، بلکه در یک
شبیه‌ساز رباتیکی دارای مدل فیزیکی، مفصل‌های گیمبال، حسگر دوربین،
سکوی پرنده و کنترلر پرواز اجرا شود.

نسخه‌های اصلی به‌کاررفته در این اجرا عبارت بودند از:

\begin{itemize}
	\item
	\EN{Ubuntu 24.04.4 LTS} در محیط \EN{WSL2}؛
	
	\item
	\EN{Gazebo Sim 8.14.0} از خانواده \EN{Harmonic}؛
	
	\item
	\EN{ArduPilot} در
	\EN{commit ff89dab}؛
	
	\item
	افزونه رسمی
	\EN{ardupilot\_gazebo}
	در
	\EN{commit 082a0fe}؛
	
	\item
	\EN{MAVProxy 1.8.74}؛
	
	\item
	\EN{GStreamer 1.24.2}
	همراه با رمزگذار
	\EN{x264 H.264}.
\end{itemize}

پلاگین‌های
\EN{ArduPilotPlugin}،
\EN{CameraZoomPlugin}،
\EN{GstCameraPlugin}
و
\EN{ParachutePlugin}
با موفقیت ساخته شدند. همچنین فایل اجرایی کنترلر پرواز
\path{build/sitl/bin/arducopter}
برای حالت
\EN{SITL}
تولید و اجرا شد.

%----------------------------------------------------------------------
\subsubsection{نقش Gazebo و ArduPilot در حلقه آزمایش}
\label{subsubsec:gazebo_ardupilot_roles}
%----------------------------------------------------------------------

در این پیاده‌سازی،
\EN{Gazebo}
نقش دوقلوی دیجیتال محیط فیزیکی را بر عهده دارد. این محیط شامل
مدل پرنده
\EN{Iris}،
مدل گیمبال، مفصل‌های چرخشی، دوربین، هدف متحرک و مدل حرکت اجزای
سامانه است. در مقابل،
\EN{ArduPilot SITL}
نرم‌افزار کنترل پرواز را اجرا می‌کند و از طریق رابط
\EN{JSON}
با مدل دینامیکی Gazebo تبادل داده دارد.

ارتباط میان دو زیرسامانه به‌صورت زیر برقرار شد:

\begin{equation}
	\begin{aligned}
		\text{\EN{ArduPilot SITL}}
		&\longrightarrow
		\text{فرمان‌ها و حالت کنترل پرواز}
		\\
		&\longrightarrow
		\text{\EN{ArduPilotPlugin}}
		\\
		&\longrightarrow
		\text{مدل \EN{Iris} در Gazebo},
	\end{aligned}
	\label{eq:ardupilot_gazebo_link}
\end{equation}

و در مسیر بازگشت، وضعیت شبیه‌سازی‌شده وسیله پرنده از Gazebo
به ArduPilot منتقل شد. دریافت پیام
\EN{JSON received}
در لاگ ArduPilot به‌عنوان شاهد برقرارشدن این اتصال ثبت شد.

\begin{keybox}[تمایز Gazebo و ArduPilot]
	Gazebo نرم‌افزار کنترل پرواز نیست؛ بلکه محیط، فیزیک، حسگرها،
	دوربین و اجزای مکانیکی را شبیه‌سازی می‌کند. ArduPilot نیز
	رندرکننده صحنه نیست؛ بلکه منطق کنترل پرواز و تله‌متری وسیله پرنده
	را اجرا می‌کند. اتصال این دو، امکان آزمون کنترلر پرواز و سامانه
	دیداری را در یک محیط فیزیکی مجازی فراهم می‌سازد.
\end{keybox}

%----------------------------------------------------------------------
\subsubsection{حلقه تصویری و کنترلی اجراشده}
\label{subsubsec:gazebo_visual_control_loop}
%----------------------------------------------------------------------

تصویر دوربین از پلاگین
\EN{GstCameraPlugin}
دریافت و به‌صورت جریان ویدئویی
\EN{H.264}
روی پورت
\EN{UDP 5600}
منتشر شد. موقعیت هدف برای تولید فرمان کنترل مستقیماً از حقیقت
مبنای شبیه‌ساز دریافت نشد؛ بلکه اندازه‌گیری کنترل‌کننده از
پیکسل‌های تصویر دوربین استخراج گردید.

حلقه اجراشده به‌صورت زیر بود:

\begin{equation}
	\boxed{
		\begin{aligned}
			\text{هدف در محیط Gazebo}
			&\longrightarrow
			\text{دوربین نصب‌شده روی گیمبال}
			\\
			&\longrightarrow
			\text{جریان تصویر \EN{H.264}}
			\\
			&\longrightarrow
			\text{اندازه‌گیری هدف از پیکسل}
			\\
			&\longrightarrow
			\text{ردیابی و تخمین حالت}
			\\
			&\longrightarrow
			\text{کنترل‌کننده \EN{Pan--Tilt}}
			\\
			&\longrightarrow
			\text{حرکت مفصل‌های گیمبال}
			\\
			&\longrightarrow
			\text{تغییر جهت دوربین و تصویر بعدی}.
		\end{aligned}
	}
	\label{eq:gazebo_pixel_closed_loop}
\end{equation}

بنابراین، تصویر فریم بعدی تابعی از وضعیت جدید گیمبال بود:

\begin{equation}
	I_{k+1}
	=
	\mathcal{R}_{\mathrm{Gazebo}}
	\left(
	\mathbf{x}^{\mathrm{target}}_{k+1},
	\mathbf{x}^{\mathrm{platform}}_{k+1},
	\mathbf{q}^{\mathrm{gimbal}}_{k+1}
	\right),
	\label{eq:gazebo_next_frame}
\end{equation}

که در آن
\(\mathcal{R}_{\mathrm{Gazebo}}\)
عملگر رندر دوربین،
\(\mathbf{x}^{\mathrm{target}}\)
حالت هدف،
\(\mathbf{x}^{\mathrm{platform}}\)
حالت سکوی حامل و
\(\mathbf{q}^{\mathrm{gimbal}}\)
زاویه‌های گیمبال هستند.

این وابستگی نشان می‌دهد که آزمایش اجراشده یک بازپخش ویدئوی مستقل
از فرمان نبوده است؛ زیرا فرمان کنترل، جهت گیمبال را تغییر داده و
حرکت گیمبال نیز مستقیماً محتوای تصویر بعدی را تحت تأثیر قرار داده
است.

%----------------------------------------------------------------------
\subsubsection{فرمان بازتولید و خروجی ویدئویی}
\label{subsubsec:gazebo_reproduction}
%----------------------------------------------------------------------

اجرای کامل آزمایش از ویندوز با فرمان زیر انجام می‌شود:

\begin{latin}
	\begin{verbatim}
		wsl.exe -d Ubuntu-24.04 -u hadi -- bash -lc \
		"cd /mnt/d/Sadoghi_r1/Research/DrHarati && \
		./scripts/run_real_gazebo_closed_loop.sh"
	\end{verbatim}
\end{latin}

ویدئوی ثبت‌شده اجرای حلقه‌بسته در مسیر زیر ذخیره شد:

\begin{latin}
	\begin{verbatim}
		results/gazebo_closed_loop/20260714_225805/videos/
		gazebo_iris_gimbal_closed_loop.mp4
	\end{verbatim}
\end{latin}

مدت ویدئوی خروجی
\EN{66.0 s}
و مدت اجرای ثبت‌شده حلقه Gazebo
\EN{63.90 s}
بود. در مجموع
\EN{660}
فریم با نرخ
\EN{10 fps}
ثبت شدند.

%----------------------------------------------------------------------
\subsubsection{نتایج اندازه‌گیری‌شده}
\label{subsubsec:gazebo_measured_results}
%----------------------------------------------------------------------

جدول~\ref{tab:gazebo_closed_loop_results}
مهم‌ترین خروجی‌های اندازه‌گیری‌شده این مرحله را نشان می‌دهد.

\begin{table}[H]
	\centering
	\caption{
		خلاصه نتایج اجرای واقعی حلقه‌بسته
		\EN{Gazebo--ArduPilot SIL}.
		تمام مقادیر از اجرای ثبت‌شده استخراج شده‌اند.
	}
	\label{tab:gazebo_closed_loop_results}
	\small
	\begin{tabularx}{0.97\textwidth}{
			>{\raggedleft\arraybackslash}p{6.2cm}
			>{\centering\arraybackslash}X
		}
		\toprule
		شاخص &
		مقدار ثبت‌شده
		\\
		\midrule
		
		نوع آزمایش &
		\EN{Gazebo--ArduPilot Closed-Loop SIL}
		\\
		
		منبع تصویر &
		\EN{Gazebo GstCameraPlugin, UDP H.264, port 5600}
		\\
		
		منبع اندازه‌گیری کنترل &
		پیکسل‌های تصویر دوربین
		\\
		
		استفاده از حقیقت مبنا در کنترل &
		خیر
		\\
		
		مدت ویدئوی ثبت‌شده &
		\EN{66.0 s}
		\\
		
		مدت حلقه Gazebo &
		\EN{63.90 s}
		\\
		
		تعداد فریم‌های ثبت‌شده &
		\EN{660}
		\\
		
		نرخ ثبت فریم &
		\EN{10 fps}
		\\
		
		میانگین خطا در حالت کنترل‌کننده خاموش &
		\EN{33.04 px}
		\\
		
		میانه خطا در حالت کنترل‌کننده روشن &
		\EN{20.44 px}
		\\
		
		میانگین خطا در بخش نهایی اجرا &
		\EN{21.47 px}
		\\
		
		بیشینه دامنه حرکت محور \EN{yaw/pan} گیمبال &
		\EN{38.59 deg}
		\\
		
		بیشینه دامنه حرکت محور \EN{pitch/tilt} گیمبال &
		\EN{35.13 deg}
		\\
		
		حالات مشاهده‌شده &
		\EN{LOCK, COAST, REACQUIRE}
		\\
		
		نشانی تله‌متری ArduPilot &
		\EN{tcp:127.0.0.1:5760}
		\\
		
		تعداد پیام‌های تله‌متری ثبت‌شده &
		\EN{29}
		\\
		
		تعداد پیام‌های \EN{Heartbeat} &
		\EN{23}
		\\
		
		وضعیت ارتباط JSON میان Gazebo و ArduPilot &
		برقرار
		\\
		
		\bottomrule
	\end{tabularx}
\end{table}

دامنه حرکت ثبت‌شده در دو محور نشان می‌دهد که کنترل‌کننده صرفاً
متغیرهای نمایشی را تغییر نداده است؛ بلکه مفصل‌های گیمبال در مدل
Gazebo حرکت قابل توجهی داشته‌اند:

\begin{equation}
	\Delta\psi_{\max}
	=
	38.59^\circ,
	\qquad
	\Delta\theta_{\max}
	=
	35.13^\circ.
	\label{eq:gazebo_gimbal_motion}
\end{equation}

توالی حالات مشاهده‌شده نیز نشان‌دهنده فعال‌بودن منطق مدیریت افت
مشاهده و بازیابی هدف بود:




\begin{equation}
	\boxed{
		\mathrm{LOCK}
		\longrightarrow
		\mathrm{COAST}
		\longrightarrow
		\mathrm{REACQUIRE}
		\longrightarrow
		\mathrm{LOCK}
	}
	\label{eq:gazebo_state_machine_sequence}
\end{equation}

\begin{keybox}[شاهد اصلی بسته‌شدن حلقه]
	در این آزمایش، هدف از تصویر دوربین Gazebo اندازه‌گیری شد، فرمان
	کنترل از اندازه‌گیری پیکسلی تولید گردید، مفصل‌های گیمبال حرکت کردند
	و تصویر بعدی در نتیجه تغییر جهت دوربین عوض شد. ازاین‌رو، ارتباط
	علی میان مشاهده، کنترل، حرکت گیمبال و تصویر بعدی در محیط استاندارد
	Gazebo برقرار شده است.
\end{keybox}

%----------------------------------------------------------------------
\subsubsection{تصاویر پیشنهادی برای مستندسازی اجرا}
\label{subsubsec:gazebo_execution_figures}
%----------------------------------------------------------------------

برای ارائه روشن‌تر در نسخه نهایی، سه فریم از ویدئوی اجرای واقعی
باید استخراج و در مسیر زیر ذخیره شوند:

\begin{latin}
	\begin{verbatim}
		figures/gazebo_closed_loop/
	\end{verbatim}
\end{latin}

نام‌های پیشنهادی فایل‌ها عبارت‌اند از:

\begin{latin}
	\begin{verbatim}
		gazebo_world_and_iris.png
		gazebo_gimbal_tracking.png
		gazebo_camera_lock.png
	\end{verbatim}
\end{latin}

کد زیر به‌گونه‌ای نوشته شده است که در صورت موجودبودن تصاویر،
آن‌ها را نمایش دهد و در صورت نبودن فایل‌ها، مانع کامپایل سند نشود.

\begin{figure}[H]
	\centering
	
	\IfFileExists{figures/gazebo_closed_loop/gazebo_world_and_iris.png}{
		\begin{minipage}[t]{0.48\textwidth}
			\centering
			\includegraphics[
			width=\linewidth,
			keepaspectratio
			]{figures/gazebo_closed_loop/gazebo_world_and_iris.png}
			
			\vspace{1mm}
			\small
			نمای محیط Gazebo و مدل پرنده \EN{Iris}
		\end{minipage}
	}{
		\begin{minipage}[t]{0.48\textwidth}
			\centering
			\fbox{
				\parbox[c][4.2cm][c]{0.90\linewidth}{
					\centering
					تصویر
					\EN{gazebo\_world\_and\_iris.png}
					هنوز افزوده نشده است.
				}
			}
		\end{minipage}
	}
	\hfill
	\IfFileExists{figures/gazebo_closed_loop/gazebo_gimbal_tracking.png}{
		\begin{minipage}[t]{0.48\textwidth}
			\centering
			\includegraphics[
			width=\linewidth,
			keepaspectratio
			]{figures/gazebo_closed_loop/gazebo_gimbal_tracking.png}
			
			\vspace{1mm}
			\small
			حرکت گیمبال در پاسخ به خطای تصویری
		\end{minipage}
	}{
		\begin{minipage}[t]{0.48\textwidth}
			\centering
			\fbox{
				\parbox[c][4.2cm][c]{0.90\linewidth}{
					\centering
					تصویر
					\EN{gazebo\_gimbal\_tracking.png}
					هنوز افزوده نشده است.
				}
			}
		\end{minipage}
	}
	
	\vspace{4mm}
	
	\IfFileExists{figures/gazebo_closed_loop/gazebo_camera_lock.png}{
		\includegraphics[
		width=0.70\textwidth,
		keepaspectratio
		]{figures/gazebo_closed_loop/gazebo_camera_lock.png}
	}{
		\fbox{
			\parbox[c][4.5cm][c]{0.66\textwidth}{
				\centering
				تصویر
				\EN{gazebo\_camera\_lock.png}
				هنوز افزوده نشده است.
			}
		}
		
		\caption{
			اجرای حلقه‌بسته در محیط
			\EN{Gazebo--ArduPilot}.
			تصاویر بالا باید به‌ترتیب نمای محیط و مدل
			\EN{Iris}، حرکت گیمبال و تصویر دوربین در حالت حفظ قفل را نشان
			دهند. ورودی کنترل از پیکسل‌های تصویر دریافت شده است و حقیقت مبنا
			در تولید فرمان استفاده نشده است.
		}
		\label{fig:gazebo_closed_loop_execution}
	\end{figure}
	
	%----------------------------------------------------------------------
	\subsubsection{تفسیر علمی نتایج و محدودیت گزارش فعلی}
	\label{subsubsec:gazebo_results_interpretation}
	%----------------------------------------------------------------------
	
	نتیجه اصلی این مرحله، کاهش یک مقدار خطا به‌تنهایی نیست؛ بلکه اثبات
	اجرایی زنجیره کامل زیر است:
	
	\begin{equation}
		\text{حسگر تصویری}
		\rightarrow
		\text{اندازه‌گیری پیکسلی}
		\rightarrow
		\text{کنترل}
		\rightarrow
		\text{حرکت مفصل}
		\rightarrow
		\text{تغییر تصویر}.
		\label{eq:gazebo_main_evidence}
	\end{equation}
	
	بااین‌حال، اعداد
	\EN{33.04 px}
	و
	\EN{20.44 px}
	به‌ترتیب میانگین و میانه هستند و نباید مستقیماً برای محاسبه درصد
	بهبود با یکدیگر مقایسه شوند. در گزارش نهایی لازم است برای هر دو
	حالت خاموش و روشن، مجموعه معیارهای یکسان شامل میانگین، میانه،
	\EN{RMSE}، صدک ۹۵ و بیشینه خطا استخراج شود.
	
	همچنین تعداد پیام‌های تله‌متری ثبت‌شده نسبت به مدت آزمایش محدود
	است. در اجرای نهایی بهتر است نرخ ثبت پیام‌های
	\EN{ATTITUDE}،
	\EN{LOCAL\_POSITION}،
	\EN{HEARTBEAT}
	و وضعیت گیمبال افزایش یابد تا تحلیل زمانی دقیق‌تری ممکن شود.
	
	\begin{warningbox}[محدودیت عنوان HIL]
		با وجود استفاده از Gazebo، ArduPilot، دوربین مجازی و گیمبال دارای
		مفصل‌های فیزیکی شبیه‌سازی‌شده، هیچ دوربین، گیمبال یا کنترلر پرواز
		فیزیکی در این مرحله وارد حلقه نشده است. بنابراین عنوان صحیح این
		نتیجه
		\EN{Gazebo--ArduPilot Closed-Loop SIL}
		است، نه
		\EN{Hardware-in-the-Loop}.
	\end{warningbox}
	
	%----------------------------------------------------------------------
	\subsubsection{سطح بلوغ حاصل‌شده}
	\label{subsubsec:gazebo_maturity_level}
	%----------------------------------------------------------------------
	
	مراحل طی‌شده را می‌توان به‌صورت زیر خلاصه کرد:
	
	\begin{equation}
		\begin{aligned}
			\text{مدل عددی کنترل}
			&\longrightarrow
			\text{دموی گرافیکی مرورگری}
			\\
			&\longrightarrow
			\text{\EN{Gazebo--ArduPilot Closed-Loop SIL}}
			\\
			&\longrightarrow
			\text{گیمبال یا دوربین فیزیکی در حلقه}.
		\end{aligned}
		\label{eq:achieved_simulation_progress}
	\end{equation}
	
	\begin{decisionbox}[نتیجه مرحله Gazebo--ArduPilot]
		در این مرحله، نسخه نرم‌افزاری حلقه‌بسته سامانه با یک زیرساخت
		استاندارد رباتیکی و کنترل پرواز اجرا شده است. محیط سه‌بعدی،
		پرنده، دوربین، گیمبال، انتقال تصویر، اندازه‌گیری پیکسلی، کنترل
		\EN{Pan--Tilt}، تغییر تصویر و تله‌متری ArduPilot همگی در یک
		اجرای واحد حضور داشته‌اند. بنابراین، بخش اصلی شبیه‌سازی نرم‌افزاری
		برای ارائه و دفاع آماده است. گام بعدی برای ارتقای سطح ادعا، ورود
		یک دوربین، گیمبال یا کنترلر پرواز فیزیکی به حلقه است.
	\end{decisionbox}








\newpage
\begin{latin}
\begin{thebibliography}{99}

\bibitem{yang2022}
J. Yang, X. Liu, J. Sun, and S. Li,
``Sampled-data robust visual servoing control for moving target tracking of an inertially stabilized platform with a measurement delay,''
\EN{Automatica}, vol. 137, art. 110105, 2022.
\href{https://doi.org/10.1016/j.automatica.2021.110105}{DOI: 10.1016/j.automatica.2021.110105}.

\bibitem{hansen2024}
J. G. Hansen and R. P. de Figueiredo,
``Active Object Detection and Tracking Using Gimbal Mechanisms for Autonomous Drone Applications,''
\EN{Drones}, vol. 8, no. 2, art. 55, 2024.
\href{https://doi.org/10.3390/drones8020055}{DOI: 10.3390/drones8020055}.

\bibitem{huynh2024}
T. Huynh et al.,
``A Study on Robust Finite-Time Visual Servoing with a Gyro-Stabilized Surveillance System,''
\EN{Actuators}, vol. 13, no. 3, art. 82, 2024.
\href{https://doi.org/10.3390/act13030082}{DOI: 10.3390/act13030082}.

\bibitem{yang2021}
L. Yang et al.,
``Image-Based Visual Servo Tracking Control of a Ground Moving Target by an Unmanned Aerial Vehicle,''
\EN{Journal of Intelligent \& Robotic Systems}, 2021.
\href{https://doi.org/10.1007/s10846-021-01425-y}{DOI: 10.1007/s10846-021-01425-y}.

\bibitem{kim2025}
J. Kim et al.,
``Vision-Based Geolocation of Moving Ground Targets Using Kalman Filtering with a Gimbal Camera on Board a UAV,''
\EN{Aerospace}, vol. 12, no. 12, art. 1065, 2025.
\href{https://doi.org/10.3390/aerospace12121065}{DOI: 10.3390/aerospace12121065}.

\bibitem{kcf}
J. F. Henriques, R. Caseiro, P. Martins, and J. Batista,
``High-Speed Tracking with Kernelized Correlation Filters,''
\EN{IEEE Transactions on Pattern Analysis and Machine Intelligence}, vol. 37, no. 3, pp. 583--596, 2015.
\href{https://doi.org/10.1109/TPAMI.2014.2345390}{DOI: 10.1109/TPAMI.2014.2345390}.

\bibitem{csrt}
A. Luke\v{z}i\v{c}, T. Voj\'{i}\v{r}, L. \v{C}ehovin, J. Matas, and M. Kristan,
``Discriminative Correlation Filter with Channel and Spatial Reliability,''
in \EN{Proceedings of CVPR}, 2017.
\href{https://doi.org/10.1109/CVPR.2017.515}{DOI: 10.1109/CVPR.2017.515}.

\bibitem{lighttrack}
B. Yan, H. Peng, K. Wu, D. Wang, J. Fu, and H. Lu,
``LightTrack: Finding Lightweight Neural Networks for Object Tracking via One-Shot Architecture Search,''
in \EN{Proceedings of CVPR}, 2021.
\href{https://doi.org/10.1109/CVPR46437.2021.01488}{DOI: 10.1109/CVPR46437.2021.01488}.

\bibitem{wescam}
L3Harris,
``WESCAM MX-10: Fully Digital High Definition Multi-Sensor Imaging System,'' public product brochure.
\url{https://www.l3harris.com/sites/default/files/2020-07/ims_eo_datasheet_WESCAM-MX-10-0401AC-Brochure.pdf}, accessed July 2026.

\bibitem{djih30}
DJI Enterprise,
``Zenmuse H30 Series Specifications,'' official technical specifications.
\url{https://enterprise.dji.com/zenmuse-h30-series/specs}, accessed July 2026.

\bibitem{gremsyt7}
Gremsy,
``Gremsy T7 Specifications,'' official technical specifications.
\url{https://gremsy.com/gremsy-t7-spec}, accessed July 2026.

\bibitem{gremsymanual}
Gremsy,
``Gremsy T7 Getting Started / Controller Update Rate,'' official manual.
\url{https://gremsy.com/gremsy-t7-manual-getting-started}, accessed July 2026.

\bibitem{basler98}
Basler AG,
``a2A2440-98g5mBAS Product Documentation,'' official camera specifications.
\url{https://docs.baslerweb.com/a2a2440-98g5mbas}, accessed July 2026.

\bibitem{mavlinkgimbal}
MAVLink,
``Gimbal Protocol v2,'' official developer guide.
\url{https://mavlink.io/en/services/gimbal_v2.html}, accessed July 2026.

\bibitem{milstd810}
U.S. Department of Defense,
``MIL-STD-810H, Method 514.8: Vibration,'' 2019.
Public copy derived from the DLA ASSIST distribution.

\bibitem{iec60068}
International Electrotechnical Commission,
``IEC 60068-2-64:2008, Environmental testing -- Test Fh: Vibration, broadband random and guidance,'' 2008.
\url{https://webstore.iec.ch/en/publication/547}.

\end{thebibliography}
\end{latin}




\end{document}
